% Options for packages loaded elsewhere
\PassOptionsToPackage{unicode}{hyperref}
\PassOptionsToPackage{hyphens}{url}
\documentclass[
]{article}
\usepackage{xcolor}
\usepackage{amsmath,amssymb}
\setcounter{secnumdepth}{-\maxdimen} % remove section numbering
\usepackage{iftex}
\ifPDFTeX
  \usepackage[T1]{fontenc}
  \usepackage[utf8]{inputenc}
  \usepackage{textcomp} % provide euro and other symbols
\else % if luatex or xetex
  \usepackage{unicode-math} % this also loads fontspec
  \defaultfontfeatures{Scale=MatchLowercase}
  \defaultfontfeatures[\rmfamily]{Ligatures=TeX,Scale=1}
\fi
\usepackage{lmodern}
\ifPDFTeX\else
  % xetex/luatex font selection
\fi
% Use upquote if available, for straight quotes in verbatim environments
\IfFileExists{upquote.sty}{\usepackage{upquote}}{}
\IfFileExists{microtype.sty}{% use microtype if available
  \usepackage[]{microtype}
  \UseMicrotypeSet[protrusion]{basicmath} % disable protrusion for tt fonts
}{}
\makeatletter
\@ifundefined{KOMAClassName}{% if non-KOMA class
  \IfFileExists{parskip.sty}{%
    \usepackage{parskip}
  }{% else
    \setlength{\parindent}{0pt}
    \setlength{\parskip}{6pt plus 2pt minus 1pt}}
}{% if KOMA class
  \KOMAoptions{parskip=half}}
\makeatother
\usepackage{longtable,booktabs,array}
\newcounter{none} % for unnumbered tables
\usepackage{calc} % for calculating minipage widths
% Correct order of tables after \paragraph or \subparagraph
\usepackage{etoolbox}
\makeatletter
\patchcmd\longtable{\par}{\if@noskipsec\mbox{}\fi\par}{}{}
\makeatother
% Allow footnotes in longtable head/foot
\IfFileExists{footnotehyper.sty}{\usepackage{footnotehyper}}{\usepackage{footnote}}
\makesavenoteenv{longtable}
\usepackage{graphicx}
\makeatletter
\newsavebox\pandoc@box
\newcommand*\pandocbounded[1]{% scales image to fit in text height/width
  \sbox\pandoc@box{#1}%
  \Gscale@div\@tempa{\textheight}{\dimexpr\ht\pandoc@box+\dp\pandoc@box\relax}%
  \Gscale@div\@tempb{\linewidth}{\wd\pandoc@box}%
  \ifdim\@tempb\p@<\@tempa\p@\let\@tempa\@tempb\fi% select the smaller of both
  \ifdim\@tempa\p@<\p@\scalebox{\@tempa}{\usebox\pandoc@box}%
  \else\usebox{\pandoc@box}%
  \fi%
}
% Set default figure placement to htbp
\def\fps@figure{htbp}
\makeatother
\setlength{\emergencystretch}{3em} % prevent overfull lines
\providecommand{\tightlist}{%
  \setlength{\itemsep}{0pt}\setlength{\parskip}{0pt}}
\usepackage{bookmark}
\IfFileExists{xurl.sty}{\usepackage{xurl}}{} % add URL line breaks if available
\urlstyle{same}
\hypersetup{
  hidelinks,
  pdfcreator={LaTeX via pandoc}}

\author{}
\date{}

\begin{document}

\subsection{Московский государственный университет имени М.В. Ломоносова}\label{ux43cux43eux441ux43aux43eux432ux441ux43aux438ux439-ux433ux43eux441ux443ux434ux430ux440ux441ux442ux432ux435ux43dux43dux44bux439-ux443ux43dux438ux432ux435ux440ux441ux438ux442ux435ux442-ux438ux43cux435ux43dux438-ux43c.ux432.-ux43bux43eux43cux43eux43dux43eux441ux43eux432ux430}

\subsection{Факультет вычислительной математики и кибернетики}\label{ux444ux430ux43aux443ux43bux44cux442ux435ux442-ux432ux44bux447ux438ux441ux43bux438ux442ux435ux43bux44cux43dux43eux439-ux43cux430ux442ux435ux43cux430ux442ux438ux43aux438-ux438-ux43aux438ux431ux435ux440ux43dux435ux442ux438ux43aux438}

\subsection{Лаборатория открытых информационных технологий}\label{ux43bux430ux431ux43eux440ux430ux442ux43eux440ux438ux44f-ux43eux442ux43aux440ux44bux442ux44bux445-ux438ux43dux444ux43eux440ux43cux430ux446ux438ux43eux43dux43dux44bux445-ux442ux435ux445ux43dux43eux43bux43eux433ux438ux439}

\includegraphics[width=4.35in,height=1.70833in,alt={A description...}]{media/media/image1.jpeg}

Выпускная квалификационная работа

«Применение методов машинного обучения для оценки финансового состояния компаний»

Выполнил студент 4 курса ВВО бакалавриата

Кобзарь А.А.

Научный руководитель:

В.В. Прусаков,

программист лаборатории

открытых информационных технологий

Москва, 2026

Содержание

Введение3

\begin{enumerate}
\def\labelenumi{\arabic{enumi}.}
\item
  Объект, предмет, цели и задачи исследования.3
\item
  Описание данных.4

  \begin{enumerate}
  \def\labelenumii{\arabic{enumii}.}
  \item
    Источники.5
  \item
    Целевая переменная.8
  \item
    Признаки.8
  \end{enumerate}
\item
  Обзор литературы11
\end{enumerate}

Глава I. Описание методов машинного обучения для анализа финансовой отчетности15

Глава II. Программная реализация и экспериментальные результаты.27

Экспериментальные результаты и выводы.39

Заключение45

Литература50

Приложение54

\begin{enumerate}
\def\labelenumi{\arabic{enumi}.}
\item
  Ссылка на программную реализацию54
\item
  Признаки, используемые в обучении моделей54
\item
  Важность признаков57
\item
  Метрики качества58
\end{enumerate}

Введение

В условиях роста объема корпоративных данных и повышения требований к качеству управленческих решений анализ финансовой отчетности становится не только инструментом ретроспективной оценки, но и основой для прогнозирования финансового состояния компаний. Традиционные методы анализа, основанные на линейных зависимостях и ограниченном наборе коэффициентов, часто недостаточно чувствительны к нелинейным эффектам, временным сдвигам и сложным взаимодействиям факторов.

\begin{quote}
1. Объект, предмет, цели и задачи исследования
\end{quote}

Объектом исследования является процесс анализа финансовой отчетности компаний в задачах финансового прогнозирования и оценки риска. Предметом исследования станут методы машинного обучения, применяемые к данным финансовой отчетности и временным финансовым рядам, а также процедуры отбора признаков и интерпретации факторного вклада в прогноз. Цель исследования -- разработать и обосновать подход к анализу финансовой отчетности на основе методов машинного обучения, обеспечивающий повышение качества прогноза и содержательную экономическую интерпретацию факторов. Задачи исследования:

\begin{enumerate}
\def\labelenumi{\arabic{enumi}.}
\item
  Проанализировать теоретические подходы к финансовому анализу отчетности и обосновать переход к ML-инструментам.
\item
  Реализовать процедуры отбора и сжатия признаков для снижения размерности и устранения избыточности факторов.
\item
  Построить и сравнить модели для задач регрессии и вероятностного прогноза.
\item
  Оценить качество моделей на данных по системе метрик точности и устойчивости.
\item
  Выполнить анализ значимости признаков после ансамблевых методов и дать экономическую интерпретацию полученных факторов.
\item
  Сформулировать практические рекомендации по применению выбранных моделей в анализе финансовой отчетности.
\end{enumerate}

\begin{quote}
2. Описание данных
\end{quote}

Финансовые данные, используемые в экономико-математическом моделировании и задачах машинного обучения, отличаются высокой неоднородностью по структуре, частоте обновления, источникам формирования и степени предварительной обработки. Для корректного построения моделей и интерпретации полученных результатов целесообразно выделять несколько ключевых типов финансовых данных.

Фундаментальные данные представляют собой количественные и качественные показатели, отражающие финансово-экономическое состояние компаний, отраслей и экономики в целом. Основным источником фундаментальных данных является регламентированная финансовая отчетность, формируемая в соответствии с национальными и международными стандартами (РСБУ, МСФО), а также макроэкономическая статистика. К фундаментальным данным относятся показатели бухгалтерского баланса, отчета о финансовых результатах и отчета о движении денежных средств, такие как активы, обязательства, выручка, прибыль, капитальные затраты и долговая нагрузка. Как правило, данные данного типа имеют низкую частоту обновления (квартальную или годовую) и публикуются с временным лагом относительно отчетного периода. Существенной особенностью фундаментальных данных является необходимость строгого учета даты фактического раскрытия информации. Использование значений до момента их официальной публикации приводит к смещению выборки и искажению результатов эмпирического анализа. Кроме того, фундаментальные данные подвержены ретроспективным корректировкам, что требует аккуратной работы с версиями данных при построении обучающих выборок и тестировании моделей.

Рыночные данные отражают процесс ценообразования и торговой активности на финансовых рынках. К данному типу относятся цены финансовых инструментов, объемы торгов, доходности, котировки спроса и предложения, открытый интерес, а также производные показатели, такие как волатильность и спрэды. Рыночные данные характеризуются высокой частотой наблюдений --- от дневных значений до тиковых данных, фиксирующих каждую сделку или изменение заявки в торговом стакане. В отличие от фундаментальных данных, рыночные данные не содержат временного лага публикации и доступны участникам рынка практически в реальном времени. С точки зрения прикладной математики и машинного обучения, рыночные данные представляют особый интерес благодаря большому объему наблюдений и возможности построения динамических и стохастических моделей. Вместе с тем высокая зашумленность и нестационарность временных рядов требуют применения специальных методов фильтрации, нормализации и регуляризации.\footnote{López de Prado M. Advances in financial machine learning --- Hoboken: Wiley, 2018. --- ISBN 978-1-119-49036-9}

В данной работе используются данные за период 2011--2024 гг. Выбор временного интервала обусловлен доступностью и сопоставимостью бухгалтерской (финансовой) отчетности российских компаний.

\begin{enumerate}
\def\labelenumi{\alph{enumi}.}
\item
  Источники
\end{enumerate}

Для решения задачи оценки финансового состояния компаний методами машинного обучения в работе используются данные, полученные из нескольких взаимодополняющих источников. Применение различных информационных ресурсов обусловлено необходимостью одновременного учета как внутренних финансовых характеристик компаний, отражаемых в бухгалтерской отчетности, так и внешних рыночных оценок, характеризующих положение компании на финансовом рынке.

ПАО «Московская биржа» является основным организатором торгов на российском фондовом рынке и официальным источником информации о рыночных характеристиках публичных компаний. В рамках данной выпускной квалификационной работы данные Московской биржи используются для формирования целевой переменной, отражающей рыночную капитализацию компаний. Информация о котировках акций, количестве выпущенных ценных бумаг, а также об изменениях в структуре капитала компаний публикуется Московской биржей на регулярной основе и характеризуется высокой степенью достоверности и актуальности. Наличие исторических данных позволяет формировать временные ряды рыночной капитализации и использовать их для анализа динамики финансового состояния компаний. Использование данных ПАО «Московская биржа» в качестве целевой переменной в моделях машинного обучения позволяет связать показатели бухгалтерской отчетности компаний с их рыночной оценкой, а также исследовать, в какой степени финансовые характеристики, отраженные в отчетности, объясняют различия в капитализации российских публичных компаний.\footnote{https://iss.moex.com/iss/reference/}

Финансовая отчетность компаний, являющаяся основным источником признаков для моделей машинного обучения, формируется на основе данных Государственного информационного ресурса бухгалтерской отчетности Федеральной налоговой службы Российской Федерации (ГИР БО). Данный ресурс содержит бухгалтерскую (финансовую) отчетность компаний, представленную в соответствии с российскими стандартами бухгалтерского учета. Государственный информационный ресурс бухгалтерской отчетности ФНС России (ГИР БО) используется в работе как основной источник исходных данных о финансовых показателях российских компаний. Ресурс содержит бухгалтерскую отчетность организаций, представленную в соответствии с российскими стандартами бухгалтерского учета (РСБУ), включая бухгалтерский баланс и отчет о финансовых результатах. До конца 2010-х годов сбор и хранение бухгалтерской отчетности осуществлялись органами государственной статистики (Росстат). Компании представляли отчетность в Росстат, и эти данные использовались преимущественно для формирования макроэкономической и отраслевой статистики. В связи с этим бухгалтерская отчетность за более ранние периоды аккумулировалась в статистических информационных системах. В дальнейшем функции по централизованному сбору бухгалтерской отчетности были переданы Федеральной налоговой службе (ФНС России), на базе которой был создан Государственный информационный ресурс бухгалтерской отчетности (ГИР БО). Начиная с отчетности за 2019 год, финансовая отчетность организаций формируется и хранится в едином электронном ресурсе ФНС. На основе отчетных показателей рассчитываются финансовые коэффициенты, характеризующие ликвидность, финансовую устойчивость, рентабельность и деловую активность компаний.\footnote{https://bo.nalog.gov.ru/} В настоящий момент все данные агрегированы и доступны в виде готового структурированного датасета на платформе Hugging Face.\footnote{\url{https://huggingface.co/datasets/irlspbru/RFSD}} Создание и подготовка этого набора данных были реализованы коллективом исследователей из Европейского университета {[}3{]}.\footnote{https://github.com/irlcode/RFSD}

Для расширения набора финансовых характеристик и использования агрегированных показателей в анализе дополнительно применяются данные специализированных финансово-аналитических ресурсов, включая платформу Financemarker, предоставляющую расчетные финансовые мультипликаторы на основе отчетности компаний.\footnote{https://financemarker.ru/api/}

Справочная информация о компаниях-эмитентах, а также дополнительные сведения, необходимые для идентификации и сопоставления данных из различных источников, получены из информационных систем «Интерфакс -- RuData», а также официальных публикаций Центрального банка Российской Федерации. Данные RuData применяются для сопоставления информации, полученной из различных источников, и формирования единого набора данных для последующего анализа. В частности, сведения о компаниях используются для корректного объединения данных бухгалтерской отчетности, рыночных показателей и финансовых мультипликаторов.\footnote{https://docs.efir-net.ru/}

В работе для расширения признакового пространства моделей машинного обучения применяются макроэкономические и банковские показатели, публикуемые Центральным банком Российской Федерации (ЦБ РФ). Эти данные доступны через официальный веб-сервис ЦБ и включают как ключевые ставки и инфляцию, так и курсы валют и металлов, показатели ликвидности и операций на финансовом рынке.\footnote{https://www.cbr.ru/DailyInfoWebServ/DailyInfo.asmx}

Совместное использование указанных источников позволяет сформировать комплексный набор данных, включающий бухгалтерскую отчетность, производные финансовые показатели, рыночные оценки и справочную информацию, что обеспечивает возможность построения и анализа моделей машинного обучения для оценки финансового состояния компаний.

\begin{enumerate}
\def\labelenumi{\alph{enumi}.}
\setcounter{enumi}{1}
\item
  Целевая переменная
\end{enumerate}

В данной работе для оценки финансового состояния компаний в качестве целевой переменной используется дневная рыночная капитализация. Этот показатель определяется так: \({MC}_{i,\ j} = \ P_{i,\ j}*{N\ }_{i,\ j}\), где MC -- рыночная капитализация, P -- рыночная цена обыкновенной акции компании, N -- количество акций в обращении, i -- идентификатор компании, j -- момент времени (в данном случае, на конец дня). Рыночная капитализация учитывает одновременно множество аспектов деятельности компании: прибыльность, устойчивость, перспективы роста и риск. В отличие от отдельных бухгалтерских показателей, она агрегирует оценки инвесторов, отражая текущую и ожидаемую стоимость компании. Дневная капитализация позволяет формировать временные ряды, что дает возможность моделировать динамику финансового состояния и выявлять тренды. Также использование капитализации как целевой переменной позволяет напрямую исследовать связь между внутренними финансовыми показателями компании и рыночной оценкой, что делает модели интерпретируемыми и экономически значимыми.

\begin{enumerate}
\def\labelenumi{\alph{enumi}.}
\setcounter{enumi}{2}
\item
  Признаки
\end{enumerate}

В рамках данной работы формируется пространство размерности в 270 признаков. Использование такого числа признаков направлено на повышение информативности входных данных и обеспечение применимости современных методов машинного обучения к задаче оценки финансового состояния компаний. Анализ охватывает выборку из 207 компаний, представленных в 11 различных секторах экономики. Отраслевая принадлежность компании включена в число признаков и используется для учета структурных различий в бизнес-моделях, уровне капиталоемкости, рентабельности и финансовых рисках, характерных для отдельных секторов.

Первую и базовую группу признаков составляет финансовая отчетность компаний, содержащая ключевые показатели их хозяйственной деятельности и финансовых результатов. Она формируется на основе бухгалтерской отчетности по РСБУ и МСФО и охватывает детализированные показатели бухгалтерского баланса, отчета о финансовых результатах и отчета о движении денежных средств. В данную группу входят сведения об организации, данные о бухгалтерском балансе, отчет о финансовых результатах, отчет об изменениях капитала, отчет о движении денежных средств, а также аудиторское заключение и пояснения к бухгалтерскому балансу и отчету о финансовых результатах. В составе финансовой отчетности можно выделить несколько укрупненных групп показателей, отражающих различные аспекты финансово-хозяйственной деятельности компании и традиционно используемых в финансовом анализе: 1) показатели ликвидности и платежеспособности формируются преимущественно на основе бухгалтерского баланса и характеризуют способность компании своевременно выполнять краткосрочные обязательства. К данной группе относятся структура оборотных активов, величина денежных средств и их эквивалентов, дебиторская задолженность, краткосрочные обязательства и соотношения между ними; 2) показатели финансовой устойчивости и структуры капитала отражают степень зависимости компании от заемных источников финансирования и устойчивость ее финансовой структуры. Рассчитываются на основе данных о собственном капитале, долгосрочных и краткосрочных обязательствах, резервах и нераспределенной прибыли, представленных в бухгалтерском балансе и отчете об изменениях капитала; 3) показатели рентабельности характеризуют эффективность использования ресурсов компании и способность генерировать прибыль. Формируются на основе отчета о финансовых результатах и включают показатели прибыли на различных уровнях (операционной, до налогообложения, чистой), а также соотношения прибыли с выручкой, активами и капиталом. 4) показатели деловой активности (оборачиваемости) отражают интенсивность использования активов и капитала в процессе хозяйственной деятельности. Основаны на данных о выручке, запасах, дебиторской и кредиторской задолженности и позволяют оценить скорость оборота основных элементов баланса; 5) показатели денежных потоков формируются на основе отчета о движении денежных средств и характеризуют способность компании генерировать денежные средства по операционной деятельности, а также направления их использования в инвестиционной и финансовой деятельности. Данная группа дополняет показатели прибыли и позволяет оценить качество финансовых результатов; 6) показатели накопления и распределения капитала отражаются в отчете об изменениях капитала и характеризуют процессы формирования собственного капитала, выплаты дивидендов, создание резервов и влияние переоценок и корректировок на финансовое положение компании.

Следующая группа признаков представляет собой финансовые коэффициенты и рыночные мультипликаторы, рассчитываемые на основе бухгалтерской отчетности, а также рыночной оценки компании. В совокупности они предназначены для нормализации абсолютных финансовых показателей и обеспечения сопоставимости компаний различного масштаба, отраслевой принадлежности и структуры капитала. В набор входят инвестиционные мультипликаторы, отражающие соотношение рыночной стоимости компании или ее предприятия с показателями, прибыли, выручки и денежных потоков. Эти коэффициенты широко используются для оценки относительной привлекательности компаний и позволяют выявлять переоцененные и недооцененные активы при межфирменном сравнении. Показатели рентабельности характеризуют эффективность использования активов, капитала и операционных ресурсов, а также структуру формирования прибыли на различных уровнях. Они отражают способность компании генерировать доходность из доступных ей ресурсов и являются ключевыми индикаторами качества бизнеса. Показатели долговой нагрузки и финансового риска описывают уровень заемного финансирования и устойчивость компании к обслуживанию долга. Данные коэффициенты позволяют оценить чувствительность финансового положения компании к изменению процентных ставок и снижению операционных доходов. Показатели ликвидности и оборотного капитала отражают краткосрочную платежеспособность и эффективность управления оборотными активами и обязательствами, дополняя анализ устойчивости финансовой структуры. Отдельную группу образуют инвестиционные и дивидендные показатели, характеризующие инвестиционную активность компании, политику распределения прибыли и баланс между реинвестированием и выплатами акционерам. Наличие каждого показателя в двух вариантах --- по РСБУ и МСФО --- позволяет учитывать различия в учетных стандартах и повышает устойчивость моделей к специфике раскрытия отчетности, обеспечивая более полное и надежное описание финансового состояния компаний.

Отдельную группу признаков составляют макроэкономические и денежно-кредитные показатели, отражающие общее состояние экономики и условия функционирования финансового рынка. В данную группу включены ключевая ставка Банка России и показатели инфляции, которые характеризуют стоимость капитала и уровень макроэкономической стабильности, а также параметры ликвидности банковского сектора и операций Банка России, отражающие объемы предоставления и изъятия ликвидности. Дополнительно используются показатели инструментов денежного рынка и операций рефинансирования, позволяющие учитывать изменения в механизмах монетарной политики. В качестве внешних макрофинансовых факторов в модель включены цены на драгоценные металлы, валютные курсы, а также показатели денежного рынка RUONIA, характеризующие краткосрочную стоимость рублевой ликвидности. Использование данной группы признаков позволяет учитывать влияние макроэкономической конъюнктуры и финансовых условий на финансовое состояние и рыночную оценку компаний.

Развернутое описание каждого используемого признака вынесено в приложение к выпускной квалификационной работе.

\begin{enumerate}
\def\labelenumi{\arabic{enumi}.}
\item
  Обзор литературы
\end{enumerate}

Разделим научную литературу, используемую в данной работе, на три группы. В первой группе рассматриваются исследования, посвящённые применению методов машинного обучения в финансах в целом. Здесь машинное обучение выступает как инструмент для улучшения прогнозирования доходностей, отбора факторов и построения инвестиционных стратегий в условиях сложных, высокоразмерных и шумных данных. Литература показывает, что ML-подходы способны выявлять нелинейные зависимости, сложные взаимодействия между признаками и структурные изменения на рынке, которые классические линейные модели часто не улавливают. Важным аспектом является также способность машинного обучения работать с большим числом факторов одновременно, снижая необходимость ручного отбора признаков. При этом отдельная подгруппа работ фокусируется на сравнении классических эконометрических методов и современных алгоритмов ML, показывая преимущества последних в предсказательной точности. Однако исследователи подчёркивают, что статистическое улучшение не всегда автоматически приводит к устойчивой экономической прибыли, потому что результаты чувствительны к качеству данных, выбору метрик валидации, рыночным ограничениям и транзакционным издержкам. Современный консенсус смещается от противопоставления методов к их комбинации: машинное обучение используется для извлечения сложного сигнала, а эконометрический аппарат --- для интерпретации, проверки устойчивости и экономической обоснованности полученных результатов. Дополнительно отмечается, что интеграция ML с эконометрикой позволяет также выявлять новые финансовые факторы и уточнять существующие модели ценообразования.\footnote{Avramov D., Cheng S., Metzker L. Machine Learning versus Economic Restrictions: Evidence from Stock Return Predictability // Working paper. --- 2023. --- URL: \url{https://ssrn.com/abstract=4481306}

  Coqueret G., Guida T. Machine learning for factor investing: Python version --- Boca Raton: CRC Press, 2023. --- ISBN 978-1-003-12159-6

  Gu S., Kelly B., Xiu D. Empirical asset pricing via machine learning // Journal of Financial Economics. --- 2020. --- Vol. 138, No. 3. --- P. 622--650. --- DOI: 10.1016/j.jfineco.2020.07.006 --- URL: \url{https://www.sciencedirect.com/science/article/pii/S0304405X20301813}

  Kelly B. T., Xiu D. Financial machine learning // Working paper. --- 2023. --- URL: \url{https://bfi.uchicago.edu/wp-content/uploads/2023/07/BFI_WP_2023-100.pdf}

  López de Prado M. Advances in financial machine learning --- Hoboken: Wiley, 2018. --- ISBN 978-1-119-49036-9

  Mai J., Zhang S., Zhao H., Pan L. Factor investment or feature selection analysis? // Mathematics. --- 2024. --- Vol. 13, No. 1: 9 --- DOI: 10.3390/math13010009 --- URL: \url{https://www.mdpi.com/2227-7390/13/1/9}

  Nagel S. Machine Learning in Asset Pricing / Princeton University Press. --- Princeton, NJ: Princeton University Press, 2021. --- 160\,p. --- ISBN 9780691218700 --- DOI: 10.23943/princeton/9780691218717 --- URL: \url{https://doi.org/10.23943/princeton/9780691218717}

  Shi C. From Econometrics to Machine Learning: Transforming Empirical Asset Pricing // Journal of Economic Surveys. --- 2026. --- Vol. 40, No. 1. --- P. 528--548 --- DOI: 10.1111/joes.70002 --- URL: \url{https://ideas.repec.org/a/bla/jecsur/v40y2026i1p528-548.html}

  Yuan, P. Multi-factor Selection and Stock Picking Strategy Based on Random Forest and LightGBM Algorithms. --- 2025. --- Vol. 1, No. 1 (Issue 6). --- DOI: 10.61173/w66yzp49. --- URL: \url{https://aaltodoc.aalto.fi/items/543c8893-91a7-4a61-b4c4-ff42f5283a3c}}

Вторая группа научной литературы посвящена использованию ансамблевых методов для прогнозирования временных рядов. Ансамбли рассматриваются как инструмент повышения робастности моделей и улучшения их переносимости на разные рыночные домены. Переход от одиночных моделей к композиционным схемам позволяет уменьшить переобучение, интегрировать разнородные признаки и повысить стабильность предсказаний. Существенным трендом является развитие вероятностного прогнозирования: ансамбли дают не только точечные предсказания, но и распределения, квантили и интервалы неопределённости, что особенно важно при принятии решений с учётом рыночного риска. В этой группе выделяется работа Duan et al. (NGBoost), которая предлагает вероятностное расширение градиентного бустинга. Метод ориентирован не на точечный прогноз, а на оценку условного распределения целевой переменной, что позволяет одновременно получать ожидаемое значение, интервалы прогноза и меру неопределенности. Первоначально NGBoost разрабатывался как универсальный инструмент probabilistic prediction, однако для финансовых задач его ценность особенно велика, поскольку он делает возможной оценку не только ожидаемой доходности или риска, но и вероятности неблагоприятных сценариев. Ансамблевые методы считаются одними из наиболее эффективных инструментов для прогнозирования, однако их практическая результативность зависит от дисциплины экспериментального дизайна, выбора метрик и компромисса между точностью, вычислительной стоимостью и интерпретируемостью. В современной литературе также отмечается тенденция к объединению ансамблей с глубоким обучением для улучшения представления временной структуры данных. Практическая реализация моделей для временных финансовых данных реализована в библиотеке sktime, которая представляет собой совместимый с scikit-learn фреймворк, объединяет типовые операции подготовки рядов, трансформации признаков, построения и сравнения моделей, процедуры backtesting/forecasting в согласованном API, что повышает воспроизводимость экспериментов и снижает риск методических ошибок при сравнении алгоритмов на отчетных и рыночных временных рядах (Löning et al.).\footnote{Davis R. A., Nielsen M. S. Modeling of time series using random forests: Theoretical developments // The Annals of Statistics. --- 2020. --- Vol. 48, No. 4. --- P. 2199--2223. --- DOI: 10.1214/19-AOS1875 --- URL: \url{https://doi.org/10.1214/19-AOS1875}

  Bryzgalova S., Pelger M., Zhu J. Forest through the Trees: Building Cross Sections of Stock Returns // Journal of Finance. --- 2025. --- Vol.\,80, No.\,5. --- P.\,2447--2506. --- DOI: 10.1111/jofi.13477 --- URL: \url{https://onlinelibrary.wiley.com/doi/10.1111/jofi.13477}

  Duan T., Avati A., Ding D. Y., Thai K. K., Basu S., Ng A., Schuler A. NGBoost: Natural Gradient Boosting for Probabilistic Prediction // Proceedings of the 36th International Conference on Machine Learning (ICML 2019). --- Long Beach, California, USA, 2019. --- P. 2690--2700.

  Guijo-Rubio D., Middlehurst M., Arcencio G., Furtado Silva D., Bagnall A. Unsupervised feature based algorithms for time series extrinsic regression // Data Mining and Knowledge Discovery. --- 2023. --- Vol. 37, No. 3. --- P. 1234--1259. --- DOI: 10.1007/s10618-023-00909-5 --- URL: \url{https://link.springer.com/article/10.1007/s10618-023-00909-5}

  Löning M., Bagnall A., Ganesh S., Kazakov V., Lines J., Király F. J. sktime: A unified interface for machine learning with time series // arXiv preprint. --- 2019. --- URL: \url{https://arxiv.org/abs/1909.07872}

  Machete, R. L. Contrasting probabilistic scoring rules. // Journal of Statistical Planning and Inference, 143(10) - 2013 -- p. 1781--1790.

  Martens, J. New insights and perspectives on the natural gradient method. // Technical report. -- 2014. - arXiv: 1412.1193.

  Nasios I., Vogklis K. Blending gradient boosted trees and neural networks for point and probabilistic forecasting of hierarchical time series // arXiv preprint. --- 2023. --- DOI: 10.48550/arXiv.2310.13029 --- URL: \url{https://arxiv.org/abs/2310.13029}

  Nespoli L., Medici V. Multivariate Boosted Trees and Applications to Forecasting and Control // Journal of Machine Learning Research. --- 2022. --- Vol. 23. --- P. 1--35 --- URL: \url{http://jmlr.org/papers/v23/21-0247.html}

  Tan C. W., Bergmeir C., Petitjean F., Webb G. I. Time Series Regression. --- 2022. --- URL: \url{https://arxiv.org/abs/2201.11945}}

Третья группа научной литературы сосредоточена на данных финансовой отчетности как источнике фундаментальной информации о будущих денежных потоках, устойчивости прибыли и риске компании. Показатели отчетности обладают прогностической силой и могут использоваться для оценки ожидаемой рыночной доходности. В настоящем исследовании учитываются работы, отражающие современный этап развития, где традиционные статистические методы дополняются алгоритмами машинного обучения для более точного выявления сигналов о будущем финансовом состоянии компаний. В качестве репрезентативного примера масштабной структурированной базы можно рассматривать Russian Financial Statements Database (Bondarkov et al.), где сформирован гармонизированный массив годовой неконсолидированной отчетности по широкому кругу компаний, включая не только подавших, но и не подавших отчетность. Часто литература ставит задачу прогнозирования банкротства, как в (Cheng et al.), однако подходы к анализу отчетности можно расширять и на прогнозирование будущей переоценки акций или других рыночных событий. Современная отчетность является многослойным источником данных: кроме числовых форм, она включает текстовые раскрытия, содержащие информацию о качестве корпоративной коммуникации, рисках и ожиданиях менеджмента. Поэтому теоретически обоснованным является объединение структурированных и неструктурированных признаков в едином аналитическом контуре.\footnote{Al-Kassar T., Saadat M., Moloi T., Masadeh A., Al-Hattab A., Jrairah T. The use of accounting and financial ratios to predict failure: The case of Jordan // Academy of Accounting and Financial Studies Journal. --- 2019. --- Vol. 23, No. 3. --- URL: \url{https://www.abacademies.org/articles/the-use-of-accounting-and-financial-ratios-to-predict-failure-the-case-of-jordan.pdf}

  Bondarkov S., Ledenev V., Skougarevskiy D. Russian financial statements database: A firm-level collection of the universe of financial statements // Sci. Data. --- 2025. --- Vol. 12, Art. 995. --- DOI: 10.1038/s41597-025-05150-1 --- URL: \url{https://www.nature.com/articles/s41597-025-05150-1}

  Chen T. K., Liao H. H., Chen G. D., Kang W. H., Lin Y. C. Bankruptcy prediction using machine learning models with the text-based communicative value of annual reports // Expert Syst. Appl. --- 2023. --- Vol. 233. --- Article 120714. --- DOI: 10.1016/j.eswa.2023.120714 --- URL: \url{https://www.sciencedirect.com/science/article/pii/S0957417423012162}

  Ou J. A., Penman S. H. Financial statement analysis and the prediction of stock returns // Journal of Accounting and Economics. --- 1989. --- Vol. 11, No. 4. --- P. 295--329. --- DOI: 10.1016/0165-4101(89)90017-7 --- URL: \url{https://www.sciencedirect.com/science/article/pii/0165410189900177}

  Piotroski J. D. Value investing: The use of historical financial statement information to separate winners from losers // Journal of Accounting Research. --- 2000. --- Vol. 38, Suppl. --- P. 1--41. --- DOI: 10.2307/2672906 --- URL: \url{https://www.jstor.org/stable/2672906}

  Волков Д. Л. Управление ценностью: показатели и модели оценки // Российский журнал менеджмента. --- 2005. --- Т. 3, № 4. --- С. 67--76

  Бухвалов А. В., Волков Д. Л. Фундаментальная ценность собственного капитала: использование в управлении компанией // Научные доклады. --- No R1--2005. --- СПб.: НИИ менеджмента СПбГУ, 2005

  Ивашковская И. В. Модели экономической прибыли и контроль создания стоимости компании: дискуссионные вопросы --- 2022. --- С. 22}

Глава I. Описание методов машинного обучения для анализа финансовой отчетности

Теоретическая основа анализа финансовой отчетности в данной работе строится на представлении отчетности как формализованной информационной модели бизнеса \textbf{{[}19; 20; 24; 25{]}}. Она одновременно отражает текущее финансовое состояние компании, динамику ее результатов и ограничения по устойчивости, поэтому может использоваться не только для ретроспективной диагностики, но и для прогнозирования будущих экономических исходов \textbf{{[}1; 5{]}}. В рамках этой логики учетные показатели интерпретируются как сигналы будущих денежных потоков, прибыльности и риска \textbf{{[}20; 21; 27{]}}. Предиктивная ценность возникает не столько из отдельных коэффициентов, сколько из их совместной структуры: сочетаний, относительных изменений во времени и взаимосвязей между блоками отчетности (ликвидность, рентабельность, долговая нагрузка, оборачиваемость, качество прибыли) \textbf{{[}9; 11; 17{]}}. Следовательно, предмет анализа --- это система признаков, а не изолированные метрики \textbf{{[}6; 15{]}}.

Второй ключевой элемент теоретической базы связан с рисковой интерпретацией отчетности. Финансовые показатели выступают прокси-переменными для вероятности неблагоприятных событий (финансовые затруднения, ухудшение кредитного профиля, снижение инвестиционной привлекательности) \textbf{{[}1; 5{]}}. Отсюда вытекает методологический переход от описательного финансового анализа к постановке задач регрессии и прогнозирования вероятностей \textbf{{[}8; 11{]}}.

Третий элемент --- стоимостная перспектива. Показатели отчетности рассматриваются как факторы формирования фундаментальной ценности и экономической прибыли \textbf{{[}25; 26; 27{]}}. Это позволяет связать бухгалтерские данные с задачами оценки стоимости компании и эффективности управления ценностью, а также корректно обосновать выбор целевых переменных в эмпирических моделях \textbf{{[}20; 21{]}}.

Переход от эконометрических моделей к методам машинного обучения в финансовом анализе связан не с отказом от теории, а с изменением масштаба и структуры данных. Классическая эконометрика эффективно работает в условиях ограниченного числа факторов и заранее заданной функциональной формы, однако в задачах анализа отчетности и ценообразования активов исследователь все чаще сталкивается с сотнями признаков, сложными взаимодействиями и нестабильными рыночными режимами \textbf{{[}11; 22{]}}. Главное ограничение традиционных линейных спецификаций проявляется в высокоразмерной среде: при росте числа признаков усиливаются мультиколлинеарность, чувствительность оценок к выборке и риск переобучения при попытке вручную усложнить модель. Одновременно предпосылка линейности становится слишком жесткой, поскольку реальные зависимости в финансовых данных часто имеют пороговый, асимметричный и нелинейный характер. В результате модели, хорошо интерпретируемые внутри выборки, нередко дают слабый вневыборочный прогноз \textbf{{[}11; 17{]}}.

Методология машинного обучения предлагает другой исследовательский акцент: от приоритетной проверки параметров фиксированной модели --- к систематическому поиску устойчивой прогностической зависимости. Это достигается за счет регуляризации, гибких нелинейных алгоритмов, процедур отбора признаков и строгой вневыборочной валидации. В такой рамке центральным критерием качества становится не только статистическая значимость коэффициентов, но и воспроизводимая точность прогноза на новых данных \textbf{{[}13; 22{]}}. При этом современный подход носит гибридный характер. Экономическая интерпретация, ограничения реализуемости стратегий и контроль за методологическими ошибками сохраняются как обязательные элементы исследования, а ML выступает инструментом повышения адаптивности моделей к сложной структуре финансовой информации. Таким образом, эволюция от эконометрики к машинному обучению представляет собой переход от «жестко специфицированных» моделей к данным-ориентированной, но теоретически осмысленной аналитике \textbf{{[}11; 13; 17{]}}.

Развитие ML-подходов в финансовом анализе связано с переходом от «классических» ансамблей к более гибким архитектурам, которые одновременно повышают точность, учитывают многомерность целевых переменных и дают оценку неопределенности прогноза. В этой логике бустинг выступает как базовый механизм последовательного улучшения модели, а гибридные и вероятностные конструкции --- как его прикладное расширение под специфику финансовых задач {[}8; 19{]}. Важное направление здесь --- многомерный бустинг. В отличие от стандартной постановки с одной целевой переменной, Multivariate Boosted Trees позволяют совместно моделировать несколько взаимосвязанных выходов и учитывать их кросс-корреляции. Формально для векторной цели \(\text{(}\mathbf{y}_{\mathbf{i}} \in \mathbb{R}^{\mathbb{K}}\text{)}\) модель имеет вид:

\[\widehat{\mathbf{y}_{i}} = F\left( x_{i} \right) \in \mathbb{R}^{\mathbb{K}},\quad\quad F_{t}(x) = F_{t - 1}(x) + \eta\,\mathbf{h}_{\mathbf{t}}(x),\]

а функция потерь учитывает совместную структуру ошибок:

\[\mathcal{L} = \sum_{i = 1}^{n}{\left( \mathbf{y}_{\mathbf{i}} - \widehat{\mathbf{y}_{i}} \right)^{\top}\Sigma^{- 1}\left( \mathbf{y}_{\mathbf{i}} - \widehat{\mathbf{y}_{i}} \right)} + \Omega(F),\ \]

где \(\Sigma\) -- матрица ковариации между компонентами целевого вектора. Для финансовой аналитики это принципиально, поскольку многие целевые показатели (доходность, риск, компоненты денежных потоков, показатели устойчивости) формируются как взаимозависимая система, и раздельное моделирование может терять существенную информацию о структуре данных {[}19{]}. Деревья решений и ансамблевые методы подходят для анализа финансовых данных, поскольку хорошо работают в условиях, типичных для отчетности и рыночных признаков: высокой размерности, нелинейных зависимостей, пороговых эффектов и сложных взаимодействий факторов. В отличие от жестко заданных линейных спецификаций, такие модели извлекают структуру данных более гибко и при этом сохраняют прикладную интерпретируемость через важности признаков и анализ вкладов отдельных переменных {[}4; 24{]}. С практической точки зрения полезно рассматривать ансамбли как семейство взаимодополняющих инструментов. Случайный лес обеспечивает устойчивый базовый прогноз за счет усреднения большого числа деревьев и снижения дисперсии, тогда как бустинговые алгоритмы (в частности LightGBM) последовательно улучшают аппроксимацию, акцентируясь на ранее недообъясненных наблюдениях. Эмпирические результаты для задач отбора акций подтверждают, что комбинация RF и LightGBM позволяет эффективно совмещать классификацию сигналов и регрессионную оценку целевых величин в едином исследовательском контуре, например, через связку \(\widehat{p_{i}} = P\left( y_{i} = 1\mid x_{i} \right),\quad\widehat{r_{i}} = \mathbb{E}\left\lbrack r_{i}\mid x_{i} \right\rbrack,\quad s_{i} = \widehat{p_{i}\widehat{r_{i}}}\), где s\textsubscript{i} используется как итоговый скоринг для ранжирования активов {[}23{]}.

Отдельное значение в современных задачах анализа данных имеют вероятностные модели, в которых результатом является не точечная оценка, а условное распределение целевой переменной. В этом случае модель задает не одно значение отклика y(x), а распределение вида:

\[Y \mid X = x \sim P_{\theta(x)},\theta(x) \in \mathbb{R}^{d},\]

то есть для каждого объекта \(x\)прогнозируется полный вероятностный закон P(y \textbar{} x), на основе которого могут быть получены как ожидаемое значение целевой переменной, так и интервальные оценки, квантили и иные характеристики неопределенности. Такой подход особенно важен в прикладных областях, где необходимо учитывать не только средний ожидаемый эффект, но и вероятность неблагоприятных сценариев. В частности, в финансовых приложениях управленческие решения принимаются не только на основании величины ожидаемой доходности или риска, но и с учетом распределения возможных исходов, включая вероятность экстремальных отклонений. Именно поэтому модели, способные строить калиброванные распределительные прогнозы, представляют особый интерес.

Одним из таких методов является NGBoost {[}8{]}. Данный подход ориентирован на построение не только точечных, но и распределительных прогнозов. В отличие от классических алгоритмов градиентного бустинга, которые в задачах регрессии обычно оценивают лишь условное среднее целевой переменной, NGBoost моделирует полное условное распределение P(y \textbar{} x) в рамках заранее выбранного параметрического семейства. Тем самым метод позволяет одновременно получать оценку математического ожидания отклика и количественную характеристику предсказательной неопределенности, что делает его особенно полезным в задачах, где требуется не только высокая точность, но и надежная оценка риска. Существенным достоинством данного подхода является модульность: пользователь может независимо выбирать базовый алгоритм, семейство вероятностных распределений и функцию качества. Одноименная библиотека ngboost реализует этот подход на практике и предоставляет удобный инструмент для вероятностного прогнозирования на табличных данных, сочетая вычислительную эффективность бустинга с возможностью получать информативные оценки неопределенности без обращения к более сложным и вычислительно затратным байесовским процедурам.

Обучение в NGBoost формулируется как задача минимизации функции потерь, построенной на основе строгого proper scoring rule:

\[\mathcal{L(}\theta) = \sum_{i = 1}^{n}{S(}P_{\theta(x_{i})},y_{i}).\]

В наиболее распространенном случае в качестве функции качества используется логарифмический score, что эквивалентно минимизации отрицательного логарифмического правдоподобия:

\[S(P_{\theta(x_{i})},y_{i}) = - \log p(y_{i} \mid x_{i};\theta).
\]

Таким образом, в отличие от стандартной регрессии, где минимизируется ошибка между истинным и предсказанным значением, здесь оптимизируется качество аппроксимации всего условного распределения. После обучения модель позволяет строить доверительные или, точнее, предсказательные интервалы вида

\[\mathbb{P}\text{ ⁣}\left( Y \in \left\lbrack q_{\alpha/2}(x),\text{ }q_{1 - \alpha/2}(x) \right\rbrack \mid X = x \right) \approx 1 - \alpha,
\]

где \(q_{\alpha/2}(x)\)и \(q_{1 - \alpha/2}(x)\)--- соответствующие условные квантили распределения. Именно наличие такой интервальной информации делает вероятностные методы особенно ценными в приложениях, где важно оценивать не только ожидаемое значение, но и степень неопределенности прогноза.

Ключевая научная новизна NGBoost состоит в том, что обновление параметров распределения выполняется с использованием натурального градиента. Если условное распределение задается вектором параметров \(\theta\), то прямое применение обычного градиентного спуска в пространстве параметров приводит к тому, что динамика обучения начинает зависеть от выбранной параметризации распределения. Иными словами, одно и то же распределение, записанное через разные параметры, может порождать различные траектории оптимизации. Кроме того, при совместном обучении нескольких параметров, например среднего и дисперсии, обычный градиент может приводить к несбалансированным обновлениям, когда один параметр изменяется слишком быстро, а другой --- слишком медленно. Для устранения этого недостатка в NGBoost используется натуральный градиент, учитывающий геометрию пространства вероятностных распределений:

\[{\widetilde{\nabla}}_{\theta}S(\theta,y) = I_{S}(\theta)^{- 1}\nabla_{\theta}S(\theta,y),\]

где \(S(\theta,y)\)--- выбранная функция качества, а \(I_{S}(\theta)\)--- риманова метрика статистического многообразия, индуцированная соответствующим scoring rule. В частном случае логарифмического score эта метрика совпадает с матрицей информации Фишера, поэтому шаг натурального градиента принимает вид

\[{\widetilde{\nabla}}_{\theta}\mathcal{L =}I(\theta)^{- 1}\nabla_{\theta}\mathcal{L}.\]

Это означает, что оптимизация осуществляется не просто в евклидовом пространстве параметров, а в геометрии самих распределений, что обеспечивает инвариантность к перепараметризации и более корректную настройку различных компонент параметрического вектора.

Алгоритмически NGBoost сочетает идеи градиентного бустинга и вероятностного моделирования. На каждой итерации вычисляются натуральные псевдоостатки для каждого наблюдения, после чего базовые модели аппроксимируют найденные направления обновления параметров распределения. Далее параметры обновляются по правилу

\[\theta^{(m)}(x) = \theta^{\left( m-1 \right)}(x) + \eta\text{ }f_{m}(x),\]

а в более общем виде с учетом масштабирования шаг может быть записан как

\[\theta_{i}^{(m)} = \theta_{i}^{\left( m-1 \right)} - \eta\text{ }\rho^{(m)}f^{(m)}(x_{i}),\]

где \(\eta\) --- скорость обучения, \(\rho^{(m)}\)--- коэффициент масштаба, определяемый одномерным поиском вдоль направления спуска, а \(f^{(m)}(x_{i})\) --- предсказание базового алгоритма на \(m\)-й итерации. Таким образом, NGBoost поэтапно строит функции, описывающие, как параметры распределения зависят от признаков. Важное отличие от обычного бустинга состоит в том, что аппроксимируется не целевая переменная как таковая, а направления изменения параметров распределения, и именно это делает возможным построение полноценных вероятностных прогнозов.

В модели NBRegressor базовый учитель не предсказывает целевую переменную \(y\)напрямую, а используется для поэтапной аппроксимации направлений спуска в пространстве параметров условного распределения. Поэтому принципиальное различие между моделями Ridge и DescisionTreeRegressor проявляется прежде всего в том, каким классом функций аппроксимируются натуральные псевдоостатки на каждом шаге. Иначе говоря, внешний алгоритм NGBoost в обоих случаях остается одним и тем же --- различается именно функциональное пространство, в которое проецируется направление натурального градиента. При использовании линейного учителя для каждого параметра распределения на итерации m строится линейная функция вида

\[f_{k}^{(m)}(x) = \beta_{0k}^{(m)} + x^{\top}\beta_{k}^{(m)},\]

где коэффициенты определяются из задачи

\[\underset{\beta_{0k},\text{ }\beta_{k}}{\min}\sum_{i = 1}^{n}\left( g_{ik}^{(m)}-\beta_{0k}-x_{i}^{\top}\beta_{k} \right)^{2} + \alpha \parallel \beta_{k} \parallel_{2}^{2}.\]

В рассматриваемом случае a = 0, поэтому регуляризующий член фактически отсутствует, и базовый учитель сводится к обычной линейной least-squares-аппроксимации псевдоостатков. Содержательно это означает, что на каждой итерации NGBoost вносит в каждый параметр распределения глобальную линейную поправку, одинаково заданную на всем пространстве признаков. Для нормального распределения это можно записать как

\[{\mu^{(m)}(x) = \mu^{\left( m-1 \right)}(x) - \eta\rho^{(m)}\left( \beta_{0,\mu}^{(m)}+x^{\top}\beta_{\mu}^{(m)} \right),
}{\log\sigma^{(m)}(x) = \log\sigma^{\left( m-1 \right)}(x) - \eta\rho^{(m)}\left( \beta_{0,\sigma}^{(m)}+x^{\top}\beta_{\sigma}^{(m)} \right).}\]

Поскольку сумма линейных функций остается линейной, итоговая зависимость каждого параметра от признаков также сохраняет линейную форму:

\[\theta_{k}(x) = a_{k} + x^{\top}b_{k}.\]

Следовательно, NGBRegressor с линейным базовым учителем реализует не нелинейный ансамбль в обычном смысле, а распределительную линейную модель, обучаемую пошагово с помощью натурального градиента. Ее основными достоинствами являются интерпретируемость, гладкость зависимости, низкая дисперсия оценок и сравнительно высокая устойчивость на данных, где связь между параметрами распределения и признаками близка к линейной. Такая конфигурация особенно уместна в ситуациях, когда предполагается, что и условное среднее, и параметры неопределенности изменяются примерно линейно по признакам. Вместе с тем ее выразительность ограничена: без предварительного конструирования признаков она не способна адекватно описывать пороговые эффекты, сложные взаимодействия и локальные изменения дисперсии. Кроме того, при a = 0 регуляризация со стороны самого базового учитель отсутствует, поэтому устойчивость модели достигается главным образом за счет малого learning\_rate, shrinkage и ранней остановки. При использовании дерева решений базовый учитель для каждого параметра принимает вид кусочно-постоянной функции

\[f_{k}^{(m)}(x) = \sum_{j = 1}^{J_{m}}c_{jk}^{(m)}\text{ }\mathbf{1}\{ x \in R_{j}^{(m)}\},\]

где \(R_{j}^{(m)}\)--- области признакового пространства, порожденные последовательностью осевых разбиений, а \(c_{jk}^{(m)}\)--- константа в соответствующем листе дерева. Такое дерево обучается приближать псевдоостатки локально: алгоритм подбирает разбиения так, чтобы уменьшить внутрисегментную квадратичную ошибку по значениям \(g_{ik}^{(m)}\). Критерий friedman\_mse специально адаптирован к задачам градиентного бустинга и выбирает такие разбиения, которые наиболее эффективно улучшают аппроксимацию направления спуска. При глубине max\_depth=3 каждое дерево остается слабым и формирует лишь ограниченное число локальных областей, однако уже способно моделировать нелинейности и взаимодействия признаков.

В контексте NGBoost это приводит к принципиально иной геометрии обновлений. Если линейный учитель вносит глобальную поправку сразу для всей выборки, то дерево вносит локальные поправки, различающиеся в разных областях пространства признаков. Для нормального распределения это означает, что модель может, например, в одной области признакового пространства увеличивать среднее и одновременно уменьшать неопределенность, а в другой --- делать противоположное. После суммирования многих неглубоких деревьев формируется сложная нелинейная поверхность для каждого параметра распределения, что особенно важно при наличии гетероскедастичности, структурных разрывов и скрытых взаимодействий между признаками. Именно поэтому деревья обычно оказываются более выразительным базовым классом для табличных данных: они допускают локально различную динамику как для среднего, так и для параметров масштаба распределения. Однако такая гибкость достигается ценой меньшей гладкости, более высокой вариативности оценок и слабой способности к экстраполяции вне области обучающих наблюдений. Кроме того, при ограниченном объеме данных или высоком уровне шума конфигурация деревьев решений потенциально сильнее подвержена переобучению, чем линейная.

Таким образом, итоговое сопоставление Tree и Ridge в NGBoost сводится к выбору между глобальной линейной и локальной нелинейной аппроксимацией натуральных градиентных обновлений. Линейный учитель приводит к глобальной, гладкой и интерпретируемой зависимости параметров распределения от признаков, а учитель на деревьях позволяет строить локальные, кусочно-постоянные и существенно нелинейные обновления. Поэтому линейная модель предпочтительна тогда, когда ожидается близкая к линейной структура зависимости и особенно важна интерпретируемость параметров, тогда как бустинг деревьев более уместен для табличных данных со сложными пороговыми эффектами, взаимодействиями и неоднородной дисперсией, где необходимо гибко моделировать не только условное среднее, но и форму предсказательного распределения. Следовательно, выбор между ними определяется компромиссом между интерпретируемостью и гибкостью: линейный учитель лучше соответствует гипотезе о линейной природе распределительных эффектов, тогда как учитель деревьев - гипотезе о существенно нелинейной организации как условного среднего, так и предсказательной неопределенности.

Следовательно, NGBoost можно рассматривать как эффективный и гибкий инструмент вероятностного прогнозирования, объединяющий сильные стороны градиентного бустинга и информационно-геометрического подхода к оптимизации. Его основными характеристиками являются модульность, возможность выбора семейства распределений и функции качества, инвариантность к перепараметризации, более устойчивая и сбалансированная оптимизация по сравнению с обычным градиентным спуском, способность моделировать гетероскедастичность и хорошая масштабируемость при умеренном увеличении вычислительной сложности. Благодаря этим свойствам NGBoost представляет собой важный шаг в развитии распределительных моделей машинного обучения и особенно полезен в задачах, где требуется совместно учитывать точность прогноза и достоверную оценку неопределенности {[}8; 14; 16{]}.

Также рассмотрим подход Time Series Forest, который показывает, что деревья могут быть адаптированы к временным рядам через извлечение информативных интервальных признаков и агрегирование по ансамблю, что повышает качество распознавания сложных временных паттернов, что необходимо корректного учет динамики. Для интервала \(I = \lbrack a,b\rbrack\) используются статистики вида:

\[\mu_{I} = \frac{1}{|I|}\sum_{t = a}^{b}x_{t},\quad\quad\sigma_{I}^{2} = \frac{1}{|I|}\sum_{t = a}^{b}\left( x_{t} - \mu_{I} \right)^{2},\quad\quad\beta_{I} = \frac{\sum_{t = a}^{b}{\left( t - \overline{t} \right)\left( x_{t} - \mu_{I} \right)}}{\sum_{t = a}^{b}\left( t - \overline{t} \right)^{2}},\]

а итоговый прогноз ансамбля записывается как

\[\widehat{y} = \frac{1}{M}\sum_{m = 1}^{M}{T_{m}(z)},\]

где z -- вектор интервальных признаков, а T\textsubscript{m} -- m-е дерево. При этом теоретическая состоятельность применения случайных лесов во временном контексте получает строгую поддержку в рамках нелинейных авторегрессионных процессов, где доказаны свойства сходимости при достаточно общих предпосылках {[}7{]}. В задачах факторного прайсинга значение лесов состоит в том, что они позволяют моделировать нелинейные эффекты характеристик и их взаимодействия без предварительно навязанной формы связи:

\[\mathbb{E}\left\lbrack r_{i,t + 1}\mid x_{i,t} \right\rbrack = g\left( x_{i,t} \right),\ \ g(x) = \sum_{k}^{}{g_{k}\left( x_{k} \right)} + \sum_{k < l}^{}{g_{kl}\left( x_{k},x_{l} \right)} + \cdots.\]

Это снижает риск спецификационной ошибки, характерный для линейных факторных моделей, и делает оценку премий за риск более чувствительной к гетерогенности эмитентов и рыночных режимов {[}4{]}. Тем самым деревья и ансамбли выступают не только как «технический» инструмент повышения точности, но и как способ уточнения экономической интерпретации факторов.

В задачах анализа финансовой отчетности качество модели в значительной степени определяется не только выбранным алгоритмом, но и тем, как сформировано факторное пространство. Отбор признаков после ансамблевых моделей целесообразно строить не как «разовую» процедуру фильтрации, а как двухуровневый контур: сначала извлекаются структурные временные признаки, затем их значимость агрегируется в экономически интерпретируемые факторные блоки {[}4{]}. После обучения модели Time Series Forest на описанных выше интервальных статистиках можно агрегировать важность признаков, которая формально записывается так:

\[\text{Imp}(f) = \frac{1}{M}\sum_{m = 1}^{M}{\sum_{s \in \mathcal{S}_{\mathcal{m}}(f)}^{}w_{m,s}}\,\Delta\mathcal{I}_{\mathcal{m},\mathcal{s}},\]

где \(\Delta\mathcal{I}_{\mathcal{m},\mathcal{s}}\) -- уменьшение неопределенности в узле s. Далее важности раскладываются по горизонту, что дает экономическую интерпретацию важности фактора на определенном лаге. В нашем случае лаг задается в 1 год, так как представлены данные ежегодной отчетности.

Для более глубокой экономической интерпретации результата после обучения ансамблевой модели можно оценить совместно, а не изолированно, вклад факторов путем декомпозиции вклада по временным режимам:

\[\widehat{r_{i,t + 1}} = \overline{r} + \sum_{k \in \mathcal{K}^{\star}}^{}{\phi_{k}\left( \mathbf{x}_{\mathbf{i},\mathbf{t}} \right)},\]

где \(\phi_{k}\) -- вклад фактора в прогноз. Тогда экономическая интерпретация строится на трех уровнях: 1) важность фактора в ансамбле, 2) знак и нелинейность его предельного эффекта, 3) устойчивость вклада по временным режимам {[}6{]}. Практическая реализации данной интерпретации будет выполнена с помощью информационного коэффициента Пирсона. Для каждого фактора и каждого момента времени он рассчитывается как кросс-секционная корреляция между оцененным вкладом отдельного фактора и последующей реализованной доходностью.

По итогам главы сформирована целостная методологическая рамка анализа финансовой отчетности: от фундаментальной интерпретации бухгалтерских показателей как носителей предиктивных сигналов до выбора прикладных ML-инструментов для задач прогноза. Теоретический контур задан работами по фундаментальному анализу, экономической прибыли и управлению стоимостью, а информационный контур --- сочетанием масштабных структурированных баз отчетности и текстовых раскрытий.

В методическом плане показано, что переход от эконометрики к ML обоснован ростом размерности признаков, нелинейностью связей и необходимостью устойчивого вневыборочного качества. В прикладном разрезе наилучший баланс дают ансамблевые и бустинговые подходы: деревья и леса эффективны для нелинейных табличных данных и временной динамики, бустинг и его расширения усиливают качество прогноза в многомерной и вероятностной. Дополнительно устойчивость моделей повышается за счет этапа отбора признаков и факторного дизайна.

Таким образом, практический вывод главы состоит в том, что для анализа финансовой отчетности целесообразно использовать комбинированный пайплайн: экономически осмысленные признаки, отбор и сжатие факторов, вероятностная и объяснимая интерпретация результата с обязательной проверкой вне выборки и учетом прикладных ограничений внедрения.

Глава II. Программная реализация и экспериментальные результаты

Вычислительный эксперимент данной работы состоит в том, чтобы определить, какие бухгалтерские показатели оказываются наиболее значимыми для прогнозирования динамики капитализации компаний и в какой степени использование бухгалтерской отчетности повышает качество прогноза по сравнению с базовой рыночной постановкой. Для достижения этой цели необходимо подобрать такую конфигурацию модели, которая одновременно дает устойчивое качество на валидации и позволяет интерпретировать вклад отдельных признаков. В рамках исследования проверяются две гипотезы: (1) добавление бухгалтерских факторов к базовому набору признаков приводит к улучшению качества прогноза; (2) модели на основе деревьев решений в среднем демонстрируют более высокую точность по сравнению с линейными моделями.

Состав используемых данных включает три основных блока: рыночные данные, бухгалтерские показатели и служебные признаки. Рыночный блок описывает ежедневное состояние инструмента на рынке и включает показатели капитализации компании, цены закрытия и объема торгов. Дополнительно в этом блоке присутствуют признаки, характеризующие торговый режим и параметры обращения бумаги, в частности идентификаторы режима торгов и сведения о выпуске инструмента. Бухгалтерский блок содержит показатели финансовой отчетности, присоединенные к рыночным наблюдениям. Он включает, с одной стороны, абсолютные значения отчетных статей, отражающих структуру активов, обязательств, капитала, выручки, прибыли и денежных потоков, а с другой стороны --- производные финансовые коэффициенты и мультипликаторы. В выборку входят показатели ликвидности, долговой нагрузки, маржинальности, рентабельности, инвестиционной эффективности и рыночной оценки компании. Существенной особенностью набора является наличие показателей как по РСБУ, так и по МСФО, что позволяет учитывать различия в учетных стандартах и расширяет информативность признакового пространства. Служебный блок включает переменные, необходимые для идентификации, упорядочивания и группировки наблюдений. К ним относятся идентификатор инструмента, торговая дата, идентификатор эмитента, наименование бумаги, сектор, тип акции и год наблюдения. Кроме того, в процессе подготовки данных формируются технические признаки, используемые для построения согласованной панели и последующего моделирования, например индикаторы неторговых дней и признаки, связанные с учетом дополнительных выпусков.

На этапе очистки и выравнивания панели данные приводились к единой ежедневной структуре по каждому инструменту. Сначала выполнялась фильтрация наблюдений по допустимым режимам торгов: основной режим, режим вариации расчетов, режимы для низколиквидных и неликвидных бумаг (функция filter\_boards). После этого обрабатывались случаи, связанные с дополнительными эмиссиями, в результате чего получился признак issue\_cumsum, представляющий собой флаг допэмиссии (функция drop\_additional\_filters). Далее наблюдения упорядочивались по идентификатору инструмента и торговой дате, после чего устранялись неоднозначности внутри одной даты (функция set\_index). Следующим шагом для каждого инструмента формировалась полная календарная последовательность дат с дневной частотой. Это было необходимо для приведения всех временных рядов к общей временной сетке и для корректного построения прогноза на фиксированном горизонте (функция fill\_days). После расширения календаря строился служебный бинарный признак is\_vacation, представляющий собой флаг отсутствия торгов в указанный день. На заключительном этапе выполнялось заполнение пропусков внутри каждого инструмента. Для этого применялось последовательное прямое и обратное распространение наблюдений по времени, что позволяло восстановить непрерывную панель признаков и избежать потери значительной части выборки из-за пропущенных значений (функция ffill\_bfill).

Далее логарифмические изменения с помощью функции add\_log\_returns, реализующей преобразование вида \(r_{t} = \ln\left( \frac{P_{t}}{P_{t - 1}} \right)\) отдельно внутри каждого инструмента. Для капитализации формируются признаки логарифмической доходности: log\_returns\_dailycapitalization\_1, представляющий собой логарифмическую доходность за предыдущий день, и log\_returns\_dailycapitalization\_30 --- логарифмическую доходность с лагом в один месяц. Использование логарифмических доходностей обусловлено их свойствами: они обеспечивают аддитивность во времени, стабилизируют дисперсию временного ряда и позволяют перейти от нестационарного уровня цен (или капитализации) к более стационарному процессу. Это делает их предпочтительными при построении моделей прогнозирования. В рассматриваемой постановке log\_returns\_dailycapitalization\_1 выступает в качестве целевой переменной, отражающей краткосрочную динамику капитализации. Признак log\_returns\_dailycapitalization\_30 используется как один из объясняющих факторов, поскольку он частично автокоррелирован с целевой переменной и может захватывать более долгосрочные эффекты и инерцию во временном ряду.

Блок финансовых признаков формируется из двух групп переменных: финансовые коэффициенты и мультипликаторы, и абсолютные статьи отчетности. На следующем этапе к признакам применяется сглаживание методом LOWESS (Locally Weighted Scatterplot Smoothing) -- методом локального сглаживания ряда. вместо одной глобальной модели на весь ряд в каждой точке строится маленькая локальная регрессия~по соседним наблюдениям, в результате чего получается сглаженная кривая, которая повторяет общий тренд и хорошо работает на нелинейных трендах, что свойственно этому типу данных, уменьшает влияние случайных локальных всплесков, делает более устойчивый долгосрочный сигнал, а важность признаков более интерпретируемыми. С другой стороны, чувствителен к выбору параметра frac, задающий гладкость итогового временного ряда. В проекте для сглаживания временных рядов используется предопределённая функция LOWESS (Locally Weighted Scatterplot Smoothing) из библиотеки statsmodels. Параметр frac, определяющий долю наблюдений, используемых при локальной аппроксимации, установлен равным 0.2. Данное значение было выбрано эмпирически на основе визуального анализа, как обеспечивающее наилучший баланс между степенью сглаживания и сохранением локальной структуры данных. В результате работы функции lowess\_smooth строятся сглаженные ряды smoothed\_*, которые позволяют ослабить влияние локального шума и выделить более устойчивую компоненту динамики бухгалтерских показателей. После сглаживания для сглаженных бухгалтерских признаков дополнительно рассчитываются их однодневные логарифмические изменения. Тем самым модель получает не только уровни бухгалтерских показателей, но и информацию об их краткосрочной динамике. В итоговом признаковом пространстве одновременно присутствуют исходные бухгалтерские признаки, их сглаженные значения и производные log\_returns для сглаженных рядов, что делает описание финансового состояния компании более насыщенным. На завершающем этапе выполняется преобразование признаков к формату, пригодному для обучения модели. Категориальные переменные кодируются с помощью OneHotEncoder, а сглаженные числовые признаки масштабируются: признаки с неотрицательной областью значений нормируются через MinMaxScaler, а признаки, принимающие как положительные, так и отрицательные значения, стандартизируются через StandardScaler. После этого формируется итоговая таблица признаков в едином числовом представлении.

В результате формировалась согласованная панель вида «инструмент--дата», пригодная для последующего обучения моделей и проведения валидации по расширяющемуся окну. В прикладной постановке анализируется, как комбинация рыночных и отчетных признаков влияет на точность предсказания целевой переменной в панельной временной структуре. Единицей наблюдения выступает пара «инструмент--дата», то есть запись с индексом \(\left( i,t \right)\), где \(i\)--- финансовый инструмент (идентификатор secid), а \(t\)--- торговая дата (tradedate). Таким образом, исходный массив имеет панельный формат и может быть представлен как множество наблюдений

\[\mathcal{D = \{(}i,t,X_{i,t},y_{i,t + 1})\},
\]

где \(X_{i,t}\)--- вектор признаков, доступных на момент \(t\), а \(y_{i,t + 1}\)--- прогнозируемое значение целевой переменной на следующем шаге. Ключом панели является комбинация (secid, tradedate), что обеспечивает однозначную идентификацию наблюдений и корректную сортировку во времени. Такая структура позволяет одновременно учитывать временную динамику внутри каждого инструмента и межинструментальные различия в едином экспериментальном протоколе. Практически панель является несбалансированной: для разных инструментов длина истории наблюдений может различаться из‑за различий в датах листинга, корпоративных событий и неполноты отдельных периодов. Это учитывается в процедуре формирования обучающих и тестовых выборок, где временной порядок сохраняется строго, а качество прогноза оценивается на сопоставимых срезах данных.

Практическая реализация вычислительного эксперимента организована в виде набора Jupyter notebook, каждый из которых соответствует отдельной конфигурации модели и признакового пространства. Ноутбуки выполняют роль сценариев управления экспериментом: они используют общий программный контур подготовки данных, формирования панели, построения матрицы признаков и вектора отклика, а также единый протокол временной валидации. Благодаря этому различия между отдельными запусками сводятся не к изменению логики предобработки, а к выбору модели, объема используемого признакового пространства и способа интерпретации результата.

Во всех сценариях эксперимента используется одна и та же общая последовательность действий: сначала загружается предварительно подготовленный набор данных в панельной форме, затем из него выделяются целевая переменная, матрица признаков и базовый уровень капитализации, необходимый для восстановления прогнозов из логарифмической формы. После этого данные передаются в процедуру валидации окна с минимальной длиной обучающего периода в пять лет. На каждом интервале модель обучается на накопленной исторической выборке, после чего строит прогноз для следующего временного интервала. Полученные значения оцениваются по единому набору метрик и логируются в MLflow, что обеспечивает воспроизводимость вычислительного эксперимента и сопоставимость результатов между отдельными конфигурациями. MLflow используется как единая среда фиксации результатов эксперимента: для каждого запуска сохраняются параметры модели, значения метрик качества по отдельным интервалам и их итоговые агрегированные значения. Благодаря пошаговому логированию метрик становится возможным отслеживание динамики качества по мере перехода от одного интервала к другому, а в интерфейсе MLflow эти результаты визуализируются в виде графиков. В качестве артефактов сохраняются обученные модели, что позволяет не только сравнивать конфигурации между собой, но и воспроизводить конкретный запуск, загружать полученную модель для последующего анализа и использовать ее в дальнейших экспериментах. Обучение реализовано в модуле train.py.

Первый сценарий представляет собой базовую конфигурацию вероятностной модели NGBoost без использования расширенного бухгалтерского пространства признаков. В этой постановке в модель подается только сокращенный набор рыночных и служебных признаков, включающий текущую капитализацию, индикатор неторговых дней, признак дополнительной эмиссии и лаговые логарифмические изменения капитализации. В базовой конфигурации используются только рыночные признаки без бухгалтерских факторов. Базовый набор включает 3 признака: 1) лаговый лог-остаток на глубине 30 торговых дней:

\[x_{i,t}^{(1)} = r_{i,t - 30},r_{i,t} = \ln C_{i,t} - \ln C_{i,t - 1},\]

где \(C_{i,t}\) --- капитализация инструмента i в день t; 2) флаг неторгового дня (выходной/праздник): бинарный индикатор, равный 1, если текущая торговая сессия является первой после разрыва в торгах, и 0 --- иначе; 3) значение дополнительной эмиссии issue\_cumsum - применяется не абсолютный объем дополнительной эмиссии, а накопленный показатель флагов, что имело место это событие на определенном интервале. Данный ноутбук используется для получения опорного качества прогноза в условиях минимального признакового описания. Его основная функция состоит в том, чтобы зафиксировать базовый уровень качества, относительно которого впоследствии можно оценивать прирост от добавления бухгалтерских переменных (ngboost\_base.ipynb).

Вторая серия запусков реализует основную нелинейную конфигурацию модели NGBoost на полном признаковом пространстве. В этой постановке в модель поступает уже не сокращенный, а полный набор признаков, включающий рыночные переменные, служебные признаки, бухгалтерские коэффициенты, абсолютные статьи отчетности, а также их сглаженные и производные представления. Тем самым данный ноутбук отвечает за проверку ключевой гипотезы исследования: улучшается ли качество прогноза при переходе от базового рыночного описания к расширенному пространству признаков, дополненному бухгалтерской отчетностью. По существу, именно этот сценарий реализует основную эмпирическую проверку вклада финансовой отчетности в задачу прогнозирования (ngboost\_main.ipynb).

Третий ноутбук реализует альтернативную конфигурацию той же вероятностной схемы NGBoost, но с линейным базовым алгоритмом Ridge вместо дерева решений. Содержательно этот запуск особенно важен, поскольку он позволяет провести контролируемое сравнение двух вариантов одной и той же надстройки NGBoost: нелинейной версии с деревьями и линейной версии с Ridge. Такое сопоставление дает возможность отделить эффект расширенного признакового пространства от эффекта самой архитектуры базового алгоритма. Если «деревянная» конфигурация показывает преимущество над Ridge-конфигурацией при прочих равных условиях, это служит аргументом в пользу того, что связь между целевой переменной и признаками действительно имеет нелинейный характер и не исчерпывается линейной аппроксимацией (ngboost\_ridge.ipynb).

Для обеспечения воспроизводимости вычислительного эксперимента во всех сравниваемых запусках использовался единый протокол обучения и оценки качества. Все модели обучались на одной и той же предварительно подготовленной панельной выборке, построенной по общей схеме предобработки, формирования признаков и выделения целевой переменной. Следовательно, различия в качестве прогноза между ноутбуками можно интерпретировать как результат изменения модельной конфигурации и признакового пространства, а не как следствие различий в методике валидации или в структуре обучающей выборки. С методологической точки зрения совокупность этих ноутбуков образует последовательную матрицу экспериментальных сценариев. Базовый запуск задает нижнюю опорную точку качества без бухгалтерских факторов. Полный нелинейный запуск показывает, что происходит при добавлении расширенного бухгалтерского пространства. Линейная конфигурация NGBoost позволяет проверить, насколько значима роль нелинейности внутри одной и той же вероятностной схемы. Наконец, простая линейная регрессия служит внешним ориентиром, позволяющим оценить, насколько использование бустинговой вероятностной модели превосходит более простые линейные подходы. Именно в таком сопоставлении ноутбуки образуют не набор разрозненных технических файлов, а целостную экспериментальную систему, предназначенную для выбора рабочей конфигурации и для последующего анализа вклада бухгалтерской отчетности в предсказание капитализации.

Оценка качества прогнозирования проводится по трем взаимодополняющим метрикам: WMAPE, F1 и IC. Использование именно этого набора позволяет оценить модель с трех разных точек зрения: по точности восстановления исходного уровня капитализации, по способности правильно определять направление изменения доходности и по степени согласованности между предсказанными и фактическими доходностями.

Метрика WMAPE (Weighted Mean Absolute Percentage Error) используется для оценки ошибки прогноза в исходном масштабе капитализации компаний. Поскольку модель строит прогноз в пространстве логарифмических доходностей, итоговые предсказания преобразуются обратно в уровень капитализации, после чего рассчитывается относительная ошибка:

\[WMAPE = \frac{\sum_{i,t}^{}{\mid C_{i,t} - {\widehat{C}}_{i,t} \mid}}{\sum_{i,t}^{}{\mid C_{i,t} \mid}},\]

где \(C_{i,t}\)--- фактическая капитализация компании i в момент времени t, а \({\widehat{C}}_{i,t}\)--- соответствующее прогнозное значение. Чем меньше значение WMAPE, тем выше точность модели. В отличие от обычной средней процентной ошибки, данная метрика учитывает масштаб наблюдений и потому особенно удобна для панельных данных, в которых компании существенно различаются по размеру капитализации. Оценка метрики WMAPE реализована функцией wmape\_score.

Метрика F1 применяется для оценки качества классификации направления логарифмической доходности. В данной постановке важен не только размер ошибки прогноза, но и способность модели правильно определять знак будущего изменения --- положительный или отрицательный. Для этого фактические и прогнозные логарифмические доходности переводятся в бинарную форму:

\[y_{i,t}^{dir} = \mathbf{1}(r_{i,t} > 0),{\widehat{y}}_{i,t}^{dir} = \mathbf{1}({\widehat{r}}_{i,t} > 0),
\]

где \(r_{i,t}\)--- фактическая логарифмическая доходность, \({\widehat{r}}_{i,t}\)--- прогнозная. После этого рассчитывается F1-метрика:

\[F1 = \frac{2 \cdot Precision \cdot Recall}{Precision + Recall}.\]

Таким образом, F1 характеризует, насколько хорошо модель распознает периоды роста и снижения капитализации. Чем выше значение метрики, тем точнее модель определяет направление движения. Оценка метрики F1 реализована функцией f1\_score.

Третьей метрикой является IC (Information Coefficient), широко используемая в финансовом моделировании для оценки связи между предсказанными и реализованными доходностями. В данной работе IC интерпретируется как коэффициент корреляции между прогнозными и фактическими логарифмическими доходностями:

\[IC = corr(r_{i,t},{\widehat{r}}_{i,t}).\]

В практической реализации эта величина рассчитывается как коэффициент корреляции Пирсона. Экономический смысл метрики состоит в том, что она показывает, насколько хорошо модель улавливает структуру зависимости между ожидаемыми и реализованными изменениями. Положительное и устойчиво высокое значение IC свидетельствует о том, что модель не только предсказывает направление движения, но и в целом корректно отражает относительную динамику доходностей. Оценка метрики IC реализована функцией ic\_score.

Совместное использование WMAPE, F1 и IC позволяет получить более полную картину качества модели. WMAPE оценивает ошибку в исходном масштабе капитализации, F1 --- правильность определения знака будущей доходности, а IC --- степень согласованности между прогнозируемыми и фактическими доходностями. Тем самым качество прогноза рассматривается одновременно с точки зрения точности уровней, направлений и общей финансовой информативности предсказания.

Для оценки качества прогнозирования в работе используется схема валидации с расширяющимся окном. В данной постановке на каждом шаге обучающая выборка последовательно увеличивается, тогда как тестирование выполняется на следующем временном интервале, не использованном при обучении. Такой подход позволяет имитировать реальный режим прогнозирования, в котором модель в каждый момент времени опирается только на накопленную историческую информацию и не использует сведения из будущего. Пусть панель наблюдений задана множеством пар «инструмент--дата» (i, t), где \(i\)--- финансовый инструмент, а \(t\)--- торговая дата. Тогда на \(k\)-м шаге валидации обучающая и тестовая выборки можно представить следующим образом:

\[\mathcal{D}_{\text{train}}^{(k)} = \{(i,t):t \leq T_{k}\},\mathcal{D}_{\text{test}}^{(k)} = \{(i,t):T_{k} < t \leq T_{k} + H\},\]

где T\textsubscript{k} обозначает последнюю дату обучающего окна, а H --- горизонт прогноза. После завершения очередного интервала граница обучающего окна сдвигается вперед, и наблюдения предыдущего тестового периода включаются в обучающую выборку следующего шага. Разбиение выполняется глобально по времени для всей панели сразу, поэтому все инструменты на каждом интервале обучаются и тестируются на одинаковых временных интервалах. Благодаря этому сохраняется сопоставимость результатов между различными моделями и конфигурациями признаков, а различия в метриках качества можно интерпретировать как следствие свойств модели и признакового пространства, а не как результат различий в схеме валидации. Использование схемы расширяющегося окна особенно важно в задачах финансового прогнозирования, поскольку она учитывает временную упорядоченность данных и позволяет избежать утечки информации из будущего. Тем самым обеспечивается корректность как статистической оценки качества, так и экономической интерпретации полученных результатов.

В текущей экспериментальной постановке начальное обучающее окно включает 4 полных года наблюдений. После этого на каждом интервале модель тестируется на следующем календарном году. Таким образом, горизонт прогноза составляет один год, а шаг расширения окна также равен одному году. Следовательно, каждая последующая обучающая выборка содержит всю информацию предыдущего шага и дополнительно включает еще один год наблюдений. Такая организация обеспечивает два преимущества. Во-первых, по мере продвижения во времени модель обучается на все более полном массиве исторических данных, что повышает устойчивость оценки параметров. Во-вторых, тестирование на каждом шаге осуществляется на одинаковом по длительности интервале, что делает значения метрик качества сопоставимыми между интервалами. Выбор годового горизонта прогноза и годового шага расширения окна согласуется с особенностями используемых данных. Рыночные переменные наблюдаются ежедневно, однако бухгалтерские показатели обновляются существенно реже и несут скорее среднесрочный, чем внутридневной сигнал. Поэтому годовой тестовый интервал позволяет корректно оценить вклад бухгалтерской отчетности в прогноз и при этом сохранить реалистичную временную структуру эксперимента. Механизм расширяющегося окна реализован в модуле cv.py.

Для интерпретации результатов в работе отдельно анализируется вклад признаков в прогноз в рамках выбранной tree-конфигурации NGBoostRegressor. Такой анализ позволяет перейти от оценки общего качества модели к содержательному выводу о том, какие именно переменные формируют предсказание и в какой степени бухгалтерская отчетность действительно участвует в объяснении динамики целевой переменной. Поскольку в данной конфигурации в качестве базового алгоритма используются деревья решений, после обучения модели становится возможным извлечение показателей feature importance. Эти значения интерпретируются как количественная мера вклада отдельного признака в снижение ошибки модели при построении ансамбля. Для повышения устойчивости вывода важности признаков анализируются не по одному запуску, а агрегируются по всем интервалам расширяющегося окна. Тем самым оценивается не случайный вклад признака на отдельном временном интервале, а его устойчивая роль в прогнозировании на всей выборке.

Формально для каждого признака \(x_{j}\)можно ввести усредненную важность

\[FI_{j} = \frac{1}{K}\sum_{k = 1}^{K}FI_{j}^{(k)},\]

где \(FI_{j}^{(k)}\)--- важность признака x\textsubscript{j} на \(k\)-м интервале, а K --- число интервалов в схеме расширяющегося окна. При необходимости важности дополнительно нормируются так, чтобы их сумма по всем признакам была равна единице:

\[{\widetilde{FI}}_{j} = \frac{FI_{j}}{\sum_{m = 1}^{p}FI_{m}},\]

где p --- число признаков в модели. Такая нормировка делает результаты более наглядными и позволяет интерпретировать их как относительный вклад отдельных переменных.

В прикладном анализе признаки целесообразно рассматривать не только по отдельности, но и по смысловым группам. В рамках данной работы выделяются три блока: рыночные признаки, служебные признаки и бухгалтерские признаки. К первому блоку относятся лаговые логарифмические изменения капитализации, ко второму --- признаки календарной структуры и дополнительной эмиссии, к третьему --- показатели бухгалтерской отчетности и их сглаженные производные. Это позволяет оценить не только лидирующие отдельные переменные, но и общий вклад каждого информационного блока в формирование прогноза.

По результатам интерпретации доминирующими, как правило, оказываются признаки, описывающие недавнюю динамику рыночной капитализации, поскольку именно они непосредственно связаны с краткосрочным движением целевой переменной. В первую очередь к таким переменным относится лаговый логарифмический остаток капитализации, отражающий инерционность временного ряда.

Для содержательной интерпретации результатов в работе целесообразно анализировать не только абсолютные значения важностей, но и структуру лидирующих признаков. Если верхние позиции занимают исключительно рыночные переменные, то вклад бухгалтерской отчетности в прогноз ограничен. Если же среди доминирующих признаков заметную долю составляют бухгалтерские показатели, это означает, что модель действительно использует фундаментальную информацию при формировании прогноза. В таком случае можно перейти от общего утверждения о полезности бухгалтерской отчетности к более конкретному выводу о том, какие именно аспекты финансового состояния компании оказываются наиболее значимыми: ликвидность, долговая нагрузка, рентабельность, инвестиционная активность или структура капитала.

Дополнительно интерпретация может быть проведена на уровне укрупненных групп признаков. Для этого суммарная важность признаков внутри каждой содержательной группы агрегируется:

\[FI_{G} = \sum_{j \in G}^{}{\widetilde{FI}}_{j},\]

где G обозначает группу признаков, например рыночные, служебные или бухгалтерские переменные. Такой подход позволяет оценить, какой именно блок данных дает основной вклад в прогноз. Если суммарная важность бухгалтерского блока оказывается существенной, это служит прямым аргументом в пользу того, что использование финансовой отчетности оправдано не только с точки зрения качества прогноза, но и с точки зрения интерпретируемости модели.

Дополнительно, наряду со встроенными показателями feature importance, для оценки вклада признаков будет использоваться метод permutation feature importance. Его применение позволяет проверить устойчивость выводов, полученных на основе внутренней структуры модели бустинга деревьев, и получить внешнюю оценку значимости признаков через непосредственное влияние на качество прогноза. В отличие от стандартных importance-оценок, основанных на вкладе признака в разбиения деревьев и снижение функции потерь внутри ансамбля, permutation-подход измеряет, насколько ухудшается качество модели при разрушении информации, содержащейся в конкретном признаке. Идея метода заключается в том, что после обучения модели на каждом интервале рассматривается тестовая или валидационная выборка, на которой последовательно для каждого признака x\textsubscript{j} случайным образом переставляются его значения между наблюдениями, тогда как все остальные признаки остаются без изменений. После такой перестановки снова рассчитывается метрика качества. Если качество существенно ухудшается, это означает, что данный признак действительно несет содержательную прогностическую информацию. Если же значение метрики меняется слабо, вклад признака в итоговый прогноз можно считать ограниченным.

Формально permutation-важность признака x\textsubscript{j} может быть задана как среднее снижение качества модели после случайной перестановки этого признака:

\[PI_{j} = S(f,X,y) - \frac{1}{R}\sum_{r = 1}^{R}{S\text{ ⁣}}\left( f,\pi_{j}^{(r)}(X),y \right),\]

где S(f, X, y) --- значение выбранной метрики качества для модели f на исходных данных, \(\pi_{j}^{(r)}(X)\) --- матрица признаков после r-й случайной перестановки значений j-го признака, а R --- число повторов перестановки. Чем больше значение PI\textsubscript{j}, тем сильнее ухудшается качество прогноза при разрушении информации в признаке и тем выше его значимость для модели. В рамках данной работы permutation-оценка важности используется как дополнение к стандартным feature importance. Такой подход позволяет сопоставить две интерпретации вклада признаков: внутреннюю, основанную на структуре ансамбля, и внешнюю, основанную на фактическом изменении качества прогноза. Особенно полезным это оказывается при анализе бухгалтерских факторов, поскольку permutation-метод показывает не только то, насколько часто признак используется моделью, но и то, приводит ли его наличие к реальному улучшению прогностической способности на вневыборочных данных. Если некоторый бухгалтерский показатель одновременно демонстрирует высокое встроенное значение feature importance и высокую permutation-важность, это можно рассматривать как более сильное свидетельство его реальной прогностической значимости. Напротив, если признак оказывается важным только по одной из двух мер, его роль требует более осторожной интерпретации. Оценка вклада признаков и визуализация реализованы в модуле features.py.

Экспериментальные результаты и вывод

В качестве базовой точки отсчета были сопоставлены две конфигурации модели NGBoostRegressor на сокращенном наборе рыночных и служебных признаков без использования бухгалтерской отчетности. Цель этого этапа состояла в том, чтобы отделить вклад самой архитектуры модели от вклада расширенного признакового пространства и определить, какая конфигурация лучше справляется с прогнозированием в минимальной постановке задачи. Результаты показали, что конфигурация с базовым алгоритмом на деревьях решений заметно превосходит линейную конфигурацию. Преимущество модели бустинга деревьев наблюдается по всем основным метрикам качества и носит устойчивый характер по интервалам. Это позволяет сделать вывод о том, что даже в базовой постановке зависимость между признаками и целевой переменной имеет выраженный нелинейный характер, который не может быть в полной мере уловлен линейной моделью. Иными словами, уже на минимальном наборе признаков использование деревьев решений внутри NGBoost дает существенный выигрыш по сравнению с линейной спецификацией.

Анализ результатов для базовой конфигурации показывает, что наибольший вклад в прогноз по-прежнему вносит лаговый рыночный признак, описывающий предшествующую динамику капитализации. Такой результат является ожидаемым, поскольку именно данный фактор наиболее близок к целевой переменной и отражает инерционность временного ряда. Сравнение модельных конфигураций при этом подтверждает преимущество нелинейной модели над линейной, что свидетельствует о наличии в данных нелинейных зависимостей и взаимодействий между признаками.

Если оставить в стороне чисто рыночные и служебные переменные, то среди бухгалтерских факторов прежде всего выделяются признаки, характеризующие рыночную оценку компании относительно фундаментальной базы, в частности мультипликаторы P/CF по МСФО, P/BV по МСФО и P/FCF по РСБУ. Экономический смысл данной группы состоит в том, что она отражает соотношение между рыночной стоимостью компании и ее способностью генерировать денежный поток, свободный денежный поток и накапливать собственный капитал. Тем самым модель улавливает не только сам уровень капитализации, но и то, каким образом рынок оценивает денежную и балансовую основу бизнеса. Это позволяет говорить о присутствии в сигнале устойчивой фундаментальной компоненты, связанной с относительной оценкой компании.

Существенную роль играют и показатели рентабельности, прежде всего ROA по МСФО, ROCE и ROE по РСБУ, а также валовая и операционная рентабельность по МСФО. Эти переменные характеризуют эффективность использования активов, собственного капитала и долгосрочно задействованных ресурсов, то есть внутреннее качество хозяйственной деятельности компании. Их значимость указывает на то, что модель учитывает не только внешнюю рыночную оценку фирмы, но и ее способность устойчиво создавать прибыль на имеющейся ресурсной базе. Хотя влияние каждого отдельного признака данной группы не является доминирующим, их присутствие среди значимых переменных делает интерпретацию прогноза существенно более содержательной.

Отдельное значение имеют признаки, связанные с долговой нагрузкой и финансовой устойчивостью, включая Debt Ratio по МСФО, Debt/Equity по РСБУ и показатель покрытия процентных расходов по МСФО. Экономический смысл этих факторов состоит в том, что они описывают структуру финансирования компании, степень ее зависимости от заемного капитала и способность обслуживать долговые обязательства за счет результатов основной деятельности. Тем самым модель учитывает не только потенциал доходности бизнеса, но и ограничения, обусловленные финансовым риском, уровнем долговой нагрузки и устойчивостью к ухудшению внешних условий.

Содержательно важны также показатели, характеризующие качество и динамику собственного капитала, включая совокупный финансовый результат периода, резервный капитал, историческую величину капитала, его увеличение в предшествующем периоде, а также переоценку внеоборотных активов. Эти переменные отражают глубину капитальной базы, наличие внутренних источников устойчивости и способы формирования собственного капитала. Их значимость показывает, что для интерпретации прогноза важны не только текущие финансовые результаты, но и накопленная способность компании поддерживать устойчивость баланса на более длительном горизонте.

Наконец, дополнительную информацию несут признаки, связанные с ликвидностью и инвестиционной политикой компании, включая финансовые вложения, сальдо денежных потоков от инвестиционных операций, а также операции по приобретению долговых ценных бумаг и предоставлению займов. Экономический смысл этих переменных заключается в том, что они характеризуют направления использования денежных ресурсов, особенности управления ликвидностью и характер инвестиционного поведения компании. Их присутствие среди значимых факторов показывает, что в модели учитывается не только прибыльность и структура капитала, но и практические механизмы распоряжения денежными средствами.

Важное уточнение состоит в том, что для полного набора признаков бухгалтерские переменные в среднем действительно уступают доминирующему лаговому рыночному фактору, а часть оценок их значимости может быть менее устойчивой. Однако для отобранного подмножества признаков, которые не ухудшают качество модели, характерен сравнительно малый разброс оценок важности. Это является существенным преимуществом, поскольку указывает на относительную стабильность соответствующих эффектов и позволяет рассматривать такие переменные как надежные носители фундаментальной информации.

Таким образом, бухгалтерские признаки в данной задаче не формируют основной прогнозный сигнал и не обеспечивают заметного прироста качества модели при простом расширении признакового пространства. Однако их включение не ухудшает предсказательную силу модели, а экономическое содержание и сравнительная устойчивость части из них позволяют глубже интерпретировать прогноз. Следовательно, бухгалтерская отчетность здесь выступает не заменой рыночной динамики, а ее фундаментальным дополнением, раскрывающим экономическую природу выявленного сигнала.

С практической точки зрения полученные результаты задают следующую логику дальнейшего анализа. Поскольку нелинейная конфигурация продемонстрировала наилучшее качество, а вклад большинства признаков вне лидирующего лагового фактора оказался относительно слабым и шумным, следующим этапом целесообразно выполнить отбор ограниченного числа наиболее сильных переменных. В данной работе для этого предполагается сформировать сокращенное пространство из 20 наиболее значимых признаков и повторно обучить модель уже на отобранном наборе. Такой подход должен одновременно повысить интерпретируемость модели и проверить, удается ли сохранить или улучшить качество прогноза при уменьшении размерности признакового пространства.

Полученные результаты показывают, что ключевым фактором качества прогноза в рассматриваемой задаче является не столько само расширение признакового пространства, сколько выбор модельной конфигурации. Наиболее устойчивый эффект связан с переходом от линейной базы к «деревянной»-конфигурации NGBoostRegressor: именно это изменение дает основной прирост качества. Напротив, добавление бухгалтерских признаков оказывает существенно более умеренное воздействие: в среднем оно приводит лишь к небольшому улучшению метрик либо практически не меняет результат по сравнению с базовой постановкой. Преимущество tree-конфигурации над линейной версией естественно интерпретируется через различие в выразительной способности моделей. Линейная база предполагает, что вклад каждого признака в прогноз является аддитивным и постоянным по всей выборке. Однако в панельных финансовых данных такая предпосылка часто оказывается слишком жесткой. Влияние признаков может зависеть от уровня капитализации, режима торгов, корпоративных событий, сочетания нескольких факторов и текущего состояния рынка. Деревья решений, напротив, позволяют учитывать пороговые эффекты, локальные нелинейности и взаимодействия между переменными, что делает их более адекватными для описания неоднородной структуры финансовых рядов. Дополнительным источником преимущества конфигурации с деревьями решений является высокая коррелированность признакового пространства. Бухгалтерские коэффициенты, сглаженные показатели и их производные часто содержат частично дублирующуюся информацию. В таких условиях линейная модель склонна распределять вклад между большим числом связанных переменных, что делает оценку менее устойчивой и снижает интерпретируемость. Деревья решений более гибко отбирают информативные разбиения и тем самым лучше работают в условиях высокой размерности и мультиколлинеарности. Наконец, важную роль играет сам характер задачи. Целевая переменная отражает краткосрочную динамику капитализации, в которой присутствуют инерция, асимметрия и режимные переходы. Для такой постановки линейная база оказывается недостаточной, тогда как нелинейная конфигурация лучше извлекает сигнал из сочетания рыночных и служебных признаков. Поэтому более высокое качество нелинейной модели является не случайным результатом отдельного запуска, а содержательным следствием свойств данных и структуры целевой переменной.

С экономической точки зрения ограниченный эффект бухгалтерских признаков выглядит вполне объяснимым. Целевая переменная в работе характеризует краткосрочную динамику капитализации, тогда как бухгалтерская отчетность по своей природе отражает более медленно меняющееся фундаментальное состояние компании. Между моментом формирования отчетного показателя и его влиянием на рыночную капитализацию существует временной лаг, а сам рыночный отклик на фундаментальную информацию может быть растянут во времени и частично перекрываться иными факторами --- ожиданиями инвесторов, новостным фоном, ликвидностью и общим режимом рынка. По этой причине в краткосрочном горизонте доминирующим источником сигнала остаются рыночные лаги, а бухгалтерские данные в среднем дают лишь ограниченное дополнительное уточнение. В то же время отсутствие сильного среднего эффекта не означает бесполезности бухгалтерской отчетности. Скорее результаты показывают, что ее вклад носит диффузный характер: информация распределена между большим числом переменных, каждая из которых по отдельности имеет сравнительно небольшой вес. Если отдельные показатели ликвидности, долговой нагрузки, рентабельности или рыночных мультипликаторов устойчиво появляются среди значимых признаков, это свидетельствует о наличии фундаментального сигнала, пусть и не доминирующего по силе. Иными словами, бухгалтерская отчетность не заменяет рыночную динамику как основной источник прогностической информации, но может уточнять прогноз и повышать его устойчивость в отдельных режимах.

Несмотря на содержательность полученных результатов, интерпретация выводов должна учитывать ряд ограничений. Во-первых, эффективный объем выборки оказывается меньше исходного массива данных, поскольку панель проходит через несколько этапов фильтрации и валидация требует достаточно длинной истории наблюдений. Это повышает качество и сопоставимость выборки, но одновременно сокращает число инструментов, доступных для анализа, и может ослаблять статистическую устойчивость выводов по отдельным признакам. Во-вторых, существенным ограничением остается качество бухгалтерских данных. Финансовая отчетность публикуется с временным лагом, может содержать пропуски, пересмотры и неоднородности, связанные с различиями между РСБУ и МСФО. Даже при корректной технической синхронизации такие признаки не всегда полностью отражают ту информацию, которой реально располагал рынок в конкретный момент времени. Поэтому оценка вклада бухгалтерских факторов неизбежно чувствительна к способу привязки отчетности ко временной шкале рыночных наблюдений. В-третьих, дополнительным ограничением является шум в оценках интерпретации. Метрики IC и permutation feature importance показывают, что за пределами нескольких лидирующих признаков ранжирование факторов остается нестабильным по интервалам. Это означает, что при большом числе переменных модель сталкивается с проблемой распыления сигнала. Именно поэтому логичным следующим шагом является отбор ограниченного числа наиболее сильных признаков и повторное обучение модели на сокращенном признаковом пространстве. Такой подход должен позволить одновременно повысить интерпретируемость результатов и проверить, удается ли сохранить достигнутый уровень качества при меньшей размерности данных.

Заключение

Выполненное исследование посвящено применению методов машинного обучения для оценки финансового состояния компаний на основе данных финансовой отчетности, рыночных показателей и сопутствующих макроэкономических факторов. Актуальность данной темы определяется тем, что в современных условиях объем доступной корпоративной информации постоянно растет, а задачи анализа финансового положения организаций все чаще выходят за рамки классических коэффициентных подходов и требуют использования более гибких инструментов обработки данных. Традиционные методы финансового анализа по-прежнему сохраняют свою методологическую значимость, однако в условиях высокой размерности признакового пространства, сложных нелинейных зависимостей и неоднородности временных рядов они нередко оказываются недостаточными для построения устойчивого вневыборочного прогноза. Именно поэтому обращение к методам машинного обучения в данной работе рассматривалось не как отказ от экономической интерпретации, а как естественное развитие инструментов анализа финансовой отчетности.

В ходе работы были последовательно решены поставленные исследовательские задачи. Сначала был проведен обзор теоретических подходов к анализу финансовой отчетности и современной литературы по применению машинного обучения в финансовых задачах. Это позволило показать, что переход от жестко заданных линейных моделей к более гибким ML-подходам обусловлен не модой на алгоритмы, а объективными свойствами самих данных. Финансовое состояние компании, отражаемое в отчетности, формируется под воздействием множества взаимосвязанных факторов, а рыночная реакция на фундаментальные показатели носит сложный, нелинейный и нередко асимметричный характер. В этой связи использование ансамблевых и вероятностных моделей представляется обоснованным и методически оправданным.

Далее была сформирована эмпирическая база исследования. В работе использовалось расширенное признаковое пространство, включающее бухгалтерские показатели компаний, рыночные признаки, служебные переменные, а также макроэкономические и банковские индикаторы. Такая постановка позволила перейти от узкого анализа отдельных коэффициентов к более полной многомерной характеристике состояния компании и условий ее функционирования. При этом в качестве целевой переменной была выбрана краткосрочная динамика капитализации, что задало прикладной фокус исследования: оценить, насколько данные финансовой отчетности способны улучшить прогноз рыночного поведения компании в сочетании с более оперативной рыночной информацией. Важно подчеркнуть, что сама по себе такая постановка уже отражает современную логику финансовой аналитики, в которой фундаментальные и рыночные данные рассматриваются не изолированно, а как взаимодополняющие источники сигнала.

Содержательная часть исследования была связана с описанием и сравнением методов машинного обучения, пригодных для решения поставленной задачи. Особое внимание было уделено вероятностному бустингу NGBoost, поскольку данный подход позволяет моделировать не только точечный прогноз, но и неопределенность предсказания. Для финансовых приложений это особенно важно: практическое решение редко строится только на ожидаемом значении показателя, гораздо чаще требуется учитывать степень уверенности модели, риск ошибки и устойчивость прогноза в меняющихся рыночных режимах. В работе были рассмотрены две конфигурации данной модели --- линейная и деревья решений, что создало удобную основу для проверки ключевой гипотезы о наличии нелинейной структуры в анализируемых данных.

Практическая реализация исследования строилась на строгой схеме временной валидации с расширяющимся окном. Использование подхода расширяющегося окна позволило обеспечить корректность вневыборочной оценки и исключить утечку информации из будущего, что имеет принципиальное значение для задач финансового прогнозирования. В отличие от случайного разбиения наблюдений, такая схема воспроизводит реальный режим применения модели, когда на каждом шаге доступны только исторические данные, а тестирование проводится на следующем временном интервале. Тем самым полученные результаты можно интерпретировать не только статистически, но и экономически, как приближенные к условиям реального использования модели в аналитической практике.

Для оценки качества моделей использовалась система взаимодополняющих метрик, отражающих разные аспекты прогноза. Метрика WMAPE характеризовала точность в исходном масштабе капитализации, F1 позволяла оценить способность модели корректно определять направление изменения показателя, а IC отражал степень согласованности между прогнозируемыми и реализованными доходностями. Такой подход представляется важным достоинством работы, поскольку качество финансового прогноза не может быть сведено к одной универсальной числовой характеристике. В прикладных задачах исследователю необходимо одновременно понимать, насколько точен прогноз по уровню, насколько надежно определяется направление движения и содержит ли модель экономически полезный ранжирующий сигнал.

Ключевым результатом работы стало установление устойчивого преимущества конфигурации NGBoostRegressor с деревьями решений над линейной конфигурацией. Сравнение на базовом наборе признаков показало, что использование деревьев решений в качестве базового учителя дает заметный выигрыш по основным метрикам качества. Этот результат имеет принципиальное значение, поскольку он подтверждает наличие выраженной нелинейной структуры в зависимости между признаками и целевой переменной. Иначе говоря, даже в сокращенной постановке задачи исследуемые данные не описываются адекватно в рамках простой линейной спецификации. Влияние признаков меняется в зависимости от контекста, режима рынка, сочетания факторов и уровня самих показателей, а значит, более гибкая нелинейная модель оказывается естественно более подходящей.

Анализ значимости признаков в базовой конфигурации показал, что наибольший вклад в прогноз вносит лаговый рыночный признак, характеризующий предшествующую динамику капитализации. Данный результат является содержательно ожидаемым и согласуется с общими представлениями о финансовых временных рядах: в краткосрочном горизонте инерция и недавнее поведение рынка часто оказываются наиболее сильным источником сигнала. Вместе с тем работа показывает, что доминирование лагового признака не отменяет роли остальных переменных, а лишь указывает на неравномерность распределения прогностической информации. Существенная часть признаков вносит более слабый и менее устойчивый вклад, что особенно отчетливо проявляется в оценках permutation feature importance и IC.

Следующий важный вывод связан с добавлением бухгалтерских признаков к базовому рыночному набору. Проведенные эксперименты показали, что включение отчетности не приводит к существенному приросту качества прогноза и в большинстве случаев лишь незначительно изменяет итоговые значения метрик по сравнению с базовой моделью. Однако не менее важно и то, что добавление части бухгалтерских переменных не ухудшает предсказательную силу модели. В рамках данной постановки это означает, что роль отчетности состоит не столько в резком усилении точности прогноза, сколько в выявлении экономической составляющей сигнала и в содержательной интерпретации факторов, связанных с фундаментальным состоянием компании. Более детальный анализ показывает, что среди бухгалтерских переменных выделяется подмножество признаков, отражающих рыночную оценку компании относительно фундаментальных показателей, рентабельность, долговую нагрузку, качество собственного капитала, а также ликвидность и особенности инвестиционной политики. Именно эти переменные формируют наиболее интерпретируемую часть фундаментального сигнала. Для полного набора бухгалтерских факторов оценки значимости в среднем действительно оказываются более шумными, чем для главного лагового рыночного признака, однако для отобранных переменных, не ухудшающих качество модели, характерен сравнительно малый разброс важности. Это является существенным преимуществом: даже при умеренной индивидуальной значимости такие признаки демонстрируют относительную устойчивость по временным интервалам и потому могут рассматриваться как надежные носители экономически содержательной информации.

С экономической точки зрения полученные результаты хорошо согласуются с природой исследуемых данных. Рынок реагирует на информацию не только через формальные показатели отчетности, но и через ожидания инвесторов, новостной фон, изменения ликвидности, общее состояние экономики и особенности конкретного рыночного режима. Между возникновением фундаментального сигнала и его полной рыночной капитализацией может существовать временной лаг, а часть этого сигнала может быть заранее встроена в цены. Поэтому отсутствие резкого прироста качества после добавления отчетности следует интерпретировать не как методическую неудачу, а как эмпирически важный результат: в краткосрочном прогнозировании капитализации влияние отчетных данных ограничено и опосредовано, тогда как для более длительных горизонтов оно потенциально может оказаться сильнее.

Практическая значимость работы заключается в том, что она позволяет выстроить реалистичную логику применения методов машинного обучения в задачах анализа финансового состояния компаний. Во-первых, результаты показывают, что в подобных постановках следует уделять особое внимание не только расширению набора признаков, но и выбору модельной архитектуры. Именно переход к нелинейной конфигурации оказался главным источником прироста качества. Во-вторых, работа демонстрирует, что интерпретация ML-модели должна строиться не только на итоговых метриках, но и на анализе структуры факторов, определяющих прогноз. Это особенно важно для задач, в которых модель используется не просто как «черный ящик», а как инструмент поддержки аналитических и управленческих решений. В-третьих, предложенный подход может быть использован при построении прикладных систем мониторинга, раннего выявления ухудшения финансового положения, инвестиционного скрининга, риск-анализа и корпоративной диагностики.

Отдельного внимания заслуживают ограничения исследования. Во-первых, эффективный объем выборки после фильтрации данных и построения временной валидации оказывается меньше исходного массива, что потенциально снижает статистическую устойчивость выводов по отдельным признакам. Во-вторых, существенное влияние оказывает качество бухгалтерской информации: отчетность может содержать пропуски, пересмотры, разрывы и методологическую неоднородность, в том числе связанную с различиями между РСБУ и МСФО. В-третьих, сама интерпретация значимости признаков остается чувствительной к шуму и к конкретным временным интервалам: за пределами нескольких лидирующих факторов ранжирование переменных может меняться, что указывает на эффект распыления сигнала в высокоразмерном пространстве. Все это требует осторожности при формулировании универсальных выводов и подтверждает необходимость дальнейшего совершенствования модели.

В качестве перспектив дальнейшего исследования наиболее логичным представляется переход к сокращенному признаковому пространству, сформированному из наиболее информативных переменных. Такой шаг должен одновременно решить две задачи: повысить интерпретируемость модели и проверить, удастся ли сохранить или улучшить качество прогноза при меньшей размерности. Кроме того, перспективным направлением является более точная синхронизация бухгалтерской отчетности с рыночными данными, учет задержек публикации, а также расширение постановки на другие горизонты прогнозирования. Не исключено, что вклад фундаментальных показателей будет заметно выше в среднесрочном и долгосрочном анализе, чем в задаче краткосрочного прогнозирования капитализации. Дополнительный потенциал связан и с использованием альтернативных источников данных --- текстовой корпоративной информации, новостного фона, сигналов из отчетов менеджмента и других неструктурированных признаков.

Итогом работы стало не только сравнение конкретных моделей, но и более общий методический вывод: эффективный анализ финансового состояния компании в современных условиях требует сочетания экономической интерпретации, качественной подготовки данных и гибких ML-инструментов, способных работать с высокоразмерной и неоднородной информацией. В этом смысле поставленная цель исследования достигнута: разработан и обоснован подход к анализу финансовой отчетности на основе методов машинного обучения, выполнено сравнение модельных конфигураций, оценено качество прогноза, исследован вклад различных групп признаков и сформулированы направления дальнейшего развития модели.

Литература

\begin{enumerate}
\def\labelenumi{\arabic{enumi}.}
\item
  Al-Kassar T., Saadat M., Moloi T., Masadeh A., Al-Hattab A., Jrairah T. The use of accounting and financial ratios to predict failure: The case of Jordan // Academy of Accounting and Financial Studies Journal. --- 2019. --- Vol. 23, No. 3. --- URL: \url{https://www.abacademies.org/articles/the-use-of-accounting-and-financial-ratios-to-predict-failure-the-case-of-jordan.pdf}
\item
  Avramov D., Cheng S., Metzker L. Machine Learning versus Economic Restrictions: Evidence from Stock Return Predictability // Working paper. --- 2023. --- URL: \url{https://ssrn.com/abstract=4481306}
\item
  Bondarkov S., Ledenev V., Skougarevskiy D. Russian financial statements database: A firm-level collection of the universe of financial statements // Scientific Data. --- 2025. --- Vol. 12. --- Art. 995. --- DOI: 10.1038/s41597-025-05150-1 --- URL: \url{https://www.nature.com/articles/s41597-025-05150-1}
\item
  Bryzgalova S., Pelger M., Zhu J. Forest through the Trees: Building Cross Sections of Stock Returns // Journal of Finance. --- 2025. --- Vol. 80, No. 5. --- P. 2447--2506. --- DOI: 10.1111/jofi.13477 --- URL: \url{https://onlinelibrary.wiley.com/doi/10.1111/jofi.13477}
\item
  Chen T. K., Liao H. H., Chen G. D., Kang W. H., Lin Y. C. Bankruptcy prediction using machine learning models with the text-based communicative value of annual reports // Expert Systems with Applications. --- 2023. --- Vol. 233. --- Art. 120714. --- DOI: 10.1016/j.eswa.2023.120714 --- URL: \url{https://www.sciencedirect.com/science/article/pii/S0957417423012162}
\item
  Coqueret G., Guida T. Machine Learning for Factor Investing: Python Version. --- Boca Raton: CRC Press, 2023. --- ISBN 978-1-003-12159-6
\item
  Davis R. A., Nielsen M. S. Modeling of time series using random forests: Theoretical developments // Annals of Statistics. --- 2020. --- Vol. 48, No. 4. --- P. 2199--2223. --- DOI: 10.1214/19-AOS1875 --- URL: \url{https://doi.org/10.1214/19-AOS1875}
\item
  Duan T., Avati A., Ding D. Y., Thai K. K., Basu S., Ng A., Schuler A. NGBoost: Natural Gradient Boosting for Probabilistic Prediction // Proceedings of the 36th International Conference on Machine Learning (ICML 2019). --- Long Beach, California, USA, 2019. --- P. 2690--2700.
\item
  Gu S., Kelly B., Xiu D. Empirical asset pricing via machine learning // Journal of Financial Economics. --- 2020. --- Vol. 138, No. 3. --- P. 622--650. --- DOI: 10.1016/j.jfineco.2020.07.006 --- URL: \url{https://www.sciencedirect.com/science/article/pii/S0304405X20301813}
\item
  Guijo-Rubio D., Middlehurst M., Arcencio G., Furtado Silva D., Bagnall A. Unsupervised feature based algorithms for time series extrinsic regression // Data Mining and Knowledge Discovery. --- 2023. --- Vol. 37, No. 3. --- P. 1234--1259. --- DOI: 10.1007/s10618-023-00909-5 --- URL: \url{https://link.springer.com/article/10.1007/s10618-023-00909-5}
\item
  Kelly B. T., Xiu D. Financial machine learning // Working paper. --- 2023. --- URL: \url{https://bfi.uchicago.edu/wp-content/uploads/2023/07/BFI_WP_2023-100.pdf}
\item
  Löning M., Bagnall A., Ganesh S., Kazakov V., Lines J., Király F. J. sktime: A unified interface for machine learning with time series // arXiv preprint. --- 2019. --- URL: \url{https://arxiv.org/abs/1909.07872}
\item
  López de Prado M. Advances in Financial Machine Learning. --- Hoboken: Wiley, 2018. --- ISBN 978-1-119-49036-9
\item
  Machete, R. L. Contrasting probabilistic scoring rules. // Journal of Statistical Planning and Inference, 143(10) - 2013 -- p. 1781--1790.
\item
  Mai J., Zhang S., Zhao H., Pan L. Factor investment or feature selection analysis? // Mathematics. --- 2024. --- Vol. 13, No. 1: 9 --- DOI: 10.3390/math13010009 --- URL: \url{https://www.mdpi.com/2227-7390/13/1/9}
\item
  Martens, J. New insights and perspectives on the natural gradient method. // Technical report. -- 2014. - arXiv: 1412.1193.
\item
  Nagel S. Machine Learning in Asset Pricing. --- Princeton, NJ: Princeton University Press, 2021. --- 160\,p. --- ISBN 9780691218700 --- DOI: 10.23943/princeton/9780691218717 --- URL: \url{https://doi.org/10.23943/princeton/9780691218717}
\item
  Nasios I., Vogklis K. Blending gradient boosted trees and neural networks for point and probabilistic forecasting of hierarchical time series // arXiv preprint. --- 2023. --- DOI: 10.48550/arXiv.2310.13029 --- URL: \url{https://arxiv.org/abs/2310.13029}
\item
  Nespoli L., Medici V. Multivariate Boosted Trees and Applications to Forecasting and Control // Journal of Machine Learning Research. --- 2022. --- Vol. 23. --- P. 1--35 --- URL: \url{http://jmlr.org/papers/v23/21-0247.html}
\item
  Ou J. A., Penman S. H. Financial statement analysis and the prediction of stock returns // Journal of Accounting and Economics. --- 1989. --- Vol. 11, No. 4. --- P. 295--329. --- DOI: 10.1016/0165-4101(89)90017-7 --- URL: \url{https://www.sciencedirect.com/science/article/pii/0165410189900177}
\item
  Piotroski J. D. Value investing: The use of historical financial statement information to separate winners from losers // Journal of Accounting Research. --- 2000. --- Vol. 38, Suppl. --- P. 1--41. --- DOI: 10.2307/2672906 --- URL: \url{https://www.jstor.org/stable/2672906}
\item
  Shi C. From Econometrics to Machine Learning: Transforming Empirical Asset Pricing // Journal of Economic Surveys. --- 2026. --- Vol. 40, No. 1. --- P. 528--548 --- DOI: 10.1111/joes.70002 --- URL: \url{https://ideas.repec.org/a/bla/jecsur/v40y2026i1p528-548.html}
\item
  Tan C. W., Bergmeir C., Petitjean F., Webb G. I. Time Series Regression. --- 2022. --- URL: \url{https://arxiv.org/abs/2201.11945}
\item
  Yuan P. Multi-factor Selection and Stock Picking Strategy Based on Random Forest and LightGBM Algorithms. --- 2025. --- Vol. 1, No. 1 (Issue 6). --- DOI: 10.61173/w66yzp49. --- URL: \url{https://aaltodoc.aalto.fi/items/543c8893-91a7-4a61-b4c4-ff42f5283a3c}
\item
  Волков Д. Л. Управление ценностью: показатели и модели оценки // Российский журнал менеджмента. --- 2005. --- Т. 3, № 4. --- С. 67--76
\item
  Бухвалов А. В., Волков Д. Л. Фундаментальная ценность собственного капитала: использование в управлении компанией // Научные доклады. --- No R1--2005. --- СПб.: НИИ менеджмента СПбГУ, 2005
\item
  Ивашковская И. В. Модели экономической прибыли и контроль создания стоимости компании: дискуссионные вопросы. --- 2022
\end{enumerate}

Приложение

\begin{enumerate}
\def\labelenumi{\arabic{enumi}.}
\item
  Ссылка на программную реализацию: \url{https://github.com/alexkobz/vmk_thesis}
\item
  Признаки, используемые в обучении модели
\end{enumerate}

{\def\LTcaptype{none} % do not increment counter
\begin{longtable}[]{@{}
  >{\raggedright\arraybackslash}p{(\linewidth - 0\tabcolsep) * \real{0.9990}}@{}}
\toprule\noalign{}
\begin{minipage}[b]{\linewidth}\raggedright
{\def\LTcaptype{none} % do not increment counter
\begin{longtable}[]{@{}
  >{\raggedright\arraybackslash}p{(\linewidth - 4\tabcolsep) * \real{0.2967}}
  >{\raggedright\arraybackslash}p{(\linewidth - 4\tabcolsep) * \real{0.2747}}
  >{\raggedright\arraybackslash}p{(\linewidth - 4\tabcolsep) * \real{0.4062}}@{}}
\toprule\noalign{}
\begin{minipage}[b]{\linewidth}\raggedright
Столбец
\end{minipage} & \begin{minipage}[b]{\linewidth}\raggedright
Наименование
\end{minipage} & \begin{minipage}[b]{\linewidth}\raggedright
Комментарий
\end{minipage} \\
\midrule\noalign{}
\endhead
\bottomrule\noalign{}
\endlastfoot
debt\_equity\_rsbu, debt\_equity\_msfo & Debt / Equity & Отношение заемного капитала к собственному; отражает степень использования финансового рычага и структуру источников финансирования. \\
debt\_ratio\_rsbu, debt\_ratio\_msfo & Debt Ratio & Доля заемного капитала в общей сумме активов; характеризует уровень долговой нагрузки компании. \\
gross\_margin\_rsbu, gross\_margin\_msfo & Валовая рентабельность & Отношение валовой прибыли к выручке; характеризует эффективность основной производственной деятельности. \\
operation\_margin\_rsbu, operation\_margin\_msfo & Операционная рентабельность & Отношение операционной прибыли к выручке; отражает эффективность основной деятельности без учета финансовых и налоговых факторов. \\
interest\_coverage\_rsbu, interest\_coverage\_msfo & Interest Coverage Ratio & Отношение операционной прибыли к процентным расходам; отражает способность компании обслуживать долговые обязательства. \\
pbv\_rsbu, pbv\_msfo & P / BV & Отношение рыночной капитализации к балансовой стоимости капитала; используется для оценки переоцененности компании рынком. \\
pcf\_rsbu, pcf\_msfo & P / CF & Отношение рыночной капитализации к денежному потоку; отражает способность компании генерировать денежные средства. \\
pfcf\_rsbu, pfcf\_msfo & P / FCF & Отношение рыночной капитализации к свободному денежному потоку; применяется для оценки инвестиционной привлекательности. \\
roa\_rsbu, roa\_msfo & ROA & Рентабельность активов; характеризует эффективность использования совокупных активов компании. \\
roce\_rsbu, roce\_msfo & ROCE & Рентабельность задействованного капитала; отражает эффективность использования долгосрочного капитала. \\
roe\_rsbu, roe\_msfo & ROE & Рентабельность собственного капитала; характеризует доходность для акционеров. \\
line\_1240 & Финансовые вложения (за исключением денежных эквивалентов) & Краткосрочные инвестиции в ценные бумаги и займы. \\
line\_1340 & Переоценка внеоборотных активов & Результаты переоценки основных средств и иных внеоборотных активов. \\
line\_1360 & Резервный капитал & Средства, формируемые в соответствии с законодательством или уставом для покрытия убытков и иных целей. \\
line\_2500 & Совокупный финансовый результат периода & Суммарный финансовый результат, включающий чистую прибыль и прочие компоненты совокупного дохода. \\
line\_3100 & Величина капитала (2 года назад) & Совокупная величина собственного капитала на начало сравнительного периода. \\
line\_3210 & Увеличение капитала (1 год назад) & Суммарное увеличение собственного капитала за предыдущий год. \\
line\_4200 & Сальдо денежных потоков от инвестиционных операций & Чистый поток денежных средств, связанный с инвестиционной деятельностью компании; показывает, сколько денег было потрачено или получено от инвестиций. \\
line\_4223 & В связи с приобретением долговых ценных бумаг, предоставление займов & Расходы на инвестиции в облигации, займы другим организациям или физлицам. \\
\end{longtable}
}
\end{minipage} \\
\midrule\noalign{}
\endhead
\bottomrule\noalign{}
\endlastfoot
 \\
\end{longtable}
}

\begin{enumerate}
\def\labelenumi{\arabic{enumi}.}
\setcounter{enumi}{2}
\item
  Важность признаков
\end{enumerate}

\includegraphics[width=6.5in,height=3.9in]{media/media/image2.png}

\includegraphics[width=6.5in,height=3.9in]{media/media/image3.png}

\begin{enumerate}
\def\labelenumi{\arabic{enumi}.}
\setcounter{enumi}{3}
\item
  Метрики качества
\end{enumerate}

\includegraphics[width=6.5in,height=2.85278in]{media/media/image4.png}

\includegraphics[width=6.5in,height=2.85278in]{media/media/image5.png}

\includegraphics[width=6.5in,height=2.89236in]{media/media/image6.png}

\includegraphics[width=6.5in,height=3.09306in]{media/media/image7.png}

\end{document}
